\documentclass[11pt]{article}

\usepackage[margin=1in]{geometry}
\usepackage{amsmath,amssymb,amsthm}
\usepackage{mathtools}
\usepackage{enumitem}
\usepackage[hidelinks]{hyperref}
\usepackage[expansion=false]{microtype}
\usepackage{xr}


\newtheorem{theorem}{Theorem}[section]
\newtheorem{lemma}[theorem]{Lemma}
\newtheorem{proposition}[theorem]{Proposition}
\newtheorem{corollary}[theorem]{Corollary}
\newtheorem{conjecture}[theorem]{Conjecture}
\theoremstyle{definition}
\newtheorem{definition}[theorem]{Definition}
\newtheorem{hypothesis}[theorem]{Hypothesis}
\newtheorem{remark}[theorem]{Remark}

\DeclareMathOperator{\lcm}{lcm}
\newcommand{\NN}{\mathbb{N}}
\newcommand{\ZZ}{\mathbb{Z}}
\newcommand{\QQ}{\mathbb{Q}}
\newcommand{\PP}{\mathcal{P}}
\newcommand{\der}{{}'}
\newcommand{\dd}{\,\mathrm{d}}

\title{The level--$60$ computation for Erd\H{o}s Problem~\#307:\\ a certificate, a split sieve, and the pair sector}
\author{Vico Bonfioli\thanks{\texttt{vicobonfioli@gmail.com}.
  See the acknowledgments for a candid note on the use of AI assistance in this work.}}
\date{Version 1.75.0\quad 2 September 2026}

\externaldocument[K-]{erdos307-core}
\externaldocument[B-]{erdos307-obstructions}
\externaldocument[C-]{erdos307-analytic}
\begin{document}
\maketitle

\begin{abstract}
A solution of Erd\H{o}s Problem~\#307 is a two--cycle of the arithmetic derivative, and any solution
has $|P\cup Q|\ge60$. This note records the computation that decides the level--$60$ families.
The certificate of Proposition~\ref{prop:tailkill} empties $31{,}219$ of the $49{,}961$ admissible
bases by a Jacobi symbol, and trial division to $10^{6}$ leaves $7{,}713$. For the rest we give the
split sieve: writing a cycle on a family $S\cup M$ as $a=\prod(T\cup M)$, $b=\prod(S\setminus T)$,
the two cycle equations give $\prod T\mid\partial(\prod(S\setminus T))$ together with an identity
free of the tail. The divisibility is the identity read modulo $\prod T$, so it is a consequence
rather than a hypothesis, and by the anatomy lemma it is a congruence at each prime of $T$, costing
a factor $1/r$ each; $\prod_{p\in S}(1+1/p)$ subsets survive out of $2^{|S|}$. Neither statement
mentions the cofactor $A_S$, so the criterion reaches the families whose $A_S$ is composite and, the
tail being arbitrary, the pair sector, both of which the arity--one method cannot. Every level--$60$
base surviving the certificate is decided over all $T$ with $|T|\le6$: about $9\times10^{11}$ splits,
no factorisation, and no two--cycle. The heuristic model and a breeder form for certificate search
are included. The range $|T|\ge7$ is not decided and its price is stated.
\end{abstract}
\section{The certificate programme and its anatomy}\label{sec:certificate}

The barrier of Section~\ref{K-sec:barrier} of the core note bounds a solution; it does not exclude one. This section
records the programme that attempted the exclusion level by level, its completeness, and the point
at which it provably stops. Nothing here is used by the results of the core note.

\begin{remark}[Verification tier, and why the ladder stops here]\label{rem:tier60}
Theorem~\ref{K-thm:barrier} is machine--checked in Lean~4, and so is
Proposition~\ref{K-prop:close59}: the theorem \texttt{erdos307\_sixty} formalises the entire
argument; the forced--primes and element bounds, the Pythagorean plus--square, and a pruned
depth--first search over the pool whose \emph{completeness is proved, not assumed}
(\texttt{dfs\_sound}), the search itself running under the kernel; on top of the
two independent unverified programs. The level cannot be climbed further by finite means: already at $|U|=60$ the family
$U=\{p_1,\dots,p_{59}\}\cup\{q\}$ qualifies for \emph{every} prime $q>277$, and on it the two square
conditions read $x^2=Aq+\Pi_{59}$, $y^2=Bq+\Pi_{59}$ with $A,B=N_{59}\pm2\Pi_{59}$, i.e.\ the
Pell--type system $Bx^2-Ay^2=-4\Pi_{59}^2$ with $113$--digit coefficients. We verified that no
accessible modulus obstructs it (mod $32$, and all odd $\ell<4000$: the conditions pin
$q\equiv5\pmod 8$ and thin each $\ell\mid\Pi_{59}$ to about half of its residue classes and each
$\ell\nmid\Pi_{59}$ to about a quarter, but never to zero), consistent with the local--solubility
theme of Section~\ref{sec:anatomy}. The obstruction, however, lives at the primes \emph{of} $A$
and $B$, and it is detectable \emph{without factoring}: Proposition~\ref{prop:tailkill} below
closes this family for \emph{every} $q$ by a single Jacobi--symbol computation. (An exact sweep to
$q\le10^{12}$, \texttt{code/hunt/tail\_sweep.rs}, $1.25\times10^{11}$ candidates, none with
$Aq+\Pi_{59}$ even a single square, which corroborates it independently.) The tail family also exposes the self--similarity of the ladder:
writing $a=qm$, so that $m\,b=\Pi_{59}$ partitions the first $59$ primes, its two closing
conditions collapse to the exact equation
\[ D(m)\,D(b)+m^2\;=\;\Pi_{59} \]
together with the divisibility $m\mid D(b)$ (then $q=D(b)/m$, required prime); verified
exhaustively on a $12$--prime toy universe. Equivalently, the deficit $1-\sigma(m)\sigma(b)$ must
equal $m^2/\Pi_{59}$ \emph{exactly}: the circularity $\Delta_0=M_0N(1-s_{M_0}s_N)$ of
Section~\ref{K-sec:frame} of the core note in its purest form, with the deviation forced to be the perfect square
$m^2$. Closing even this one family is thus a $2^{58}$--fold exact--coincidence search of the same
type as \#307 itself (expected count ${\sim}2^{58}/\Pi_{59}\approx10^{-95}$): each rung of the
ladder past $60$ contains the whole problem. The one case with \emph{no} free large prime (a
$60$--set drawn entirely from small primes) is instead a finite search, and empty in the
minimal range: a direct both--squares enumeration of all $60$--prime sets with $T>2$ supported in
the first $70$ primes ($4{,}011{,}289$ sets; \texttt{code/hunt/close60.rs}) finds none passing even the
plus--square, so no level--$60$ solution has all its primes $\le349$. (The first $68$ primes give
$1{,}258{,}448$ sets, also empty, and reproduce independently.) By
Corollary~\ref{C-cor:kdetermined} this verdict is decisive rather than provisional: at these masses a
support passing the plus-- and minus--squares, with the parity of Proposition~\ref{C-prop:ladder}(2),
\emph{is} a two--cycle, so a survivor would have been a solution and not a candidate. The
enumeration extends with compute.
\end{remark}

\begin{proposition}[The canonical tail family is empty: a reciprocity kill]\label{prop:tailkill}
No solution of \#307 has support $U=\{p_1,\dots,p_{59}\}\cup\{q\}$, for any prime $q$.
\end{proposition}
\begin{proof}
A solution supported on $U$ forces $x^2=Aq+\Pi_{59}$ with $x=a+b$ and $A=N_{59}+2\Pi_{59}$
(Remark~\ref{rem:tier60}). Every prime $\ell\mid A$ exceeds $277$: for $\ell\le277$ one has
$\ell\mid\Pi_{59}$, so $\ell\mid A$ would force $\ell\mid N_{59}$, contradicting
$\gcd(N_{59},\Pi_{59})=1$ (rigidity). Hence $\ell\nmid\Pi_{59}$, and reduction mod $\ell$ gives
$x^2\equiv\Pi_{59}\pmod\ell$, requiring the Legendre symbol $(\Pi_{59}\mid\ell)\neq-1$. But the
Jacobi symbol satisfies
\[ \left(\tfrac{\Pi_{59}}{A}\right)\;=\;-1 \]
(a finite gcd--speed computation, script \texttt{code/tailkill.py}; no factorisation of the
$114$--digit $A$ is needed), so some $\ell\mid A$ of odd multiplicity has $(\Pi_{59}\mid\ell)=-1$.
The family is empty, for every $q$.

The companion symbol at $B=N_{59}-2\Pi_{59}$ equals $-1$ as well; necessarily: for any
admissible base $S$ (writing $D=D_S$, $N=N_S$, $A_S,B_S=N\pm2D$) one has $A_S\equiv B_S\pmod 8$
(as $4D\equiv8\bmod{16}$) and $A_S\equiv B_S\equiv N\pmod p$ for every odd $p\mid D$, while the
quadratic--reciprocity sign difference carries the exponent $\tfrac{p-1}{2}\cdot\tfrac{A_S-B_S}{2}$
with $\tfrac{A_S-B_S}{2}=2D$ even; hence $(D\mid A_S)=(D\mid B_S)$ always; one binary
certificate per tail family. Running it over all $49{,}961$ admissible bases of
Proposition~\ref{K-prop:close59} proves $31{,}219$ of the one--new--prime families at level $60$
empty ($62.5\%$, curiously near $5/8$); the remaining $18{,}742$ are Jacobi--inconclusive: their
$-1$--Legendre primes, if any, occur to an even count, visible only through a factorisation of
$A_S$. A factoring attack on those (trial division and Brent--Pollard $\rho$ over a Montgomery
core, \texttt{code/hunt/survivor\_kill.rs}; the found factors are kept as a pairwise--coprime basis, so
a kill certificate (an odd--multiplicity basis piece $P$ with $(D_S\mid P)=-1$) needs no
primality testing at all), pushed to $\rho$--budget $10^6$ with a Pollard $p-1$ stage, closes a
further $15{,}264$: in total $46{,}483$ of the $49{,}961$ one--new--prime families at level $60$
are proven empty ($93.0\%$). The $3{,}478$ still open ($3{,}444$ of them once
Proposition~\ref{prop:immunedecide} has emptied the $34$ immune ones) resist all of this for a
structural reason,
not a want of effort: each is Baillie--PSW \emph{composite} yet has no prime factor below ${\sim}
10^{12}$, so its $A_S$ (or $B_S$) is a \emph{balanced} semiprime of two ${\sim}57$--digit primes,
splitting one is factoring at cryptographic scale. Two cheap tests confirm no shortcut survives:
a product--tree batch--GCD over all $37{,}484$ values $A_S,B_S$ finds only shared primes below
$10^9$ (already found by $\rho$, so \emph{no} free kills), and Baillie--PSW leaves exactly
$34$ families in which $A_S$ and $B_S$ are both \emph{probable} primes; these are immune
\emph{conditionally on that primality} (a prime $A_S$ with $(D_S\mid A_S)=+1$ has no inert factor,
hence no congruence obstruction for any $q$). The conditionality is not pedantry and not
removable by more of the same computation: Baillie--PSW is rigorous only in its \emph{composite}
verdicts, so the $46{,}483$ kills are unconditional while the $34$ immunities are not. Upgrading them
would require primality \emph{certificates} for thirty-four $114$--digit integers (ECPP or a
Pocklington factorisation of $A_S-1$). The elementary route does not reach: immunity needs
\emph{both} $A_S$ and $B_S$ prime, and partial factorisation of $n-1$ (trial division to
$3\times10^6$ together with the fully split part) leaves $\log F/\log n$ below $1/3$ on at least one
side of every one of the $34$, so neither Pocklington ($F>\sqrt n$) nor
Brillhart--Lehmer--Selfridge ($F>n^{1/3}$) applies to any family: $0$ of $34$ are certifiable this
way (\texttt{code/immune\_certify.py}). This is now settled. Running the Adleman--Pomerance--Rumely
--Cohen--Lenstra test on all $68$ integers proves every one of them prime in $8$ seconds of
wall clock, so the $34$ immunities upgrade from \textup{[C]} to \textup{[P]} and are
\emph{unconditional} (\texttt{code/immune\_prove.gp}, PARI/GP $2.17.4$, whose \texttt{isprime} is a
proof and not a probable--prime test). The run reproduces the whole chain independently: $49{,}961$
admissible bases, $31{,}219$ Jacobi kills, and, among the $81$ bases whose $A_S$ and $B_S$ are both
prime, exactly the $34$ with $(D_S\mid A_S)=+1$; the other $47$ carry $(D_S\mid A_S)=-1$ and were
already killed. The primality is moreover backed by
\emph{independently checkable} evidence rather than a verdict: \texttt{code/immune\_certify.gp}
emits Atkin--Morain \textup{ECPP} certificates for all $68$ integers, archived at
\texttt{data/certs/immune\_ecpp.txt} ($512$\,KB) so a third
party can check them without trusting the prover. That file is a \texttt{primecertexport} dump and
is not machine--readable: PARI's \texttt{read} parses its header as a polynomial, so
\texttt{primecertisvalid} cannot consume it. The working check is
\texttt{code/immune\_cert\_verify.gp}, which extracts the $68$ asserted subjects and re--proves each
by \textup{ECPP} in $1.3$ seconds: $68/68$ prime, $0$ failures. The structural half of the campaign is now
formal: the identity $(D_S\mid A_S)=(D_S\mid B_S)$, which is what makes one binary test decide a
family rather than two, is proved in Lean~4 as \texttt{Erdos307.jacobiSym\_A\_eq\_B}
\textup{(\texttt{lean/Erdos307/Campaign.lean})}, on the three standard axioms and with no
\texttt{native\_decide}; it follows from periodicity, $A_S-B_S=4D_S$ being exactly the period of
$J(a\mid\cdot)$. The primality itself is not formalised and cannot presently be: mathlib's only
route to a large prime is \texttt{lucas\_primality}, which needs a complete factorisation of
$n-1$, and that is precisely what the $0$ of $34$ above rules out. Note the two counts answer different questions and both are correct: $81$ is the
number of prime pairs, $34$ the number of \emph{immune} families among them. A complete kill/immune split of the residual is in any case beyond feasible
computation, and the honest state of the single--tail sector is \emph{$46{,}483$ empty
\textup{[P]} and $3{,}478$ open}, the latter including the $34$ immune ones, whose primality the
\textup{APRCL} run above makes unconditional \textup{[P]}; Proposition~\ref{prop:immunedecide} then
empties those $34$, leaving $3{,}444$.
\end{proof}

\begin{proposition}[Sector decomposition of level $60$]\label{prop:sectors}
Let $U$ be a $60$--element prime support with $T(U)>2$. Then exactly one of the following holds:
\textup{(i)} $T(U)-2\ge 1/\max U$, and $U=S\cup\{q\}$ for an admissible base $S$
\textup{(}one of the $49{,}961$ of Proposition~\ref{K-prop:close59}\textup{)}; the
\emph{single--tail sector}, or \textup{(ii)} $T(U)-2<1/\max U$, and, writing $p>q$ for the two
largest elements and $W$ for the rest, $T(W\cup\{p\})<2$ and $T(W\cup\{q\})<2$; the
\emph{pair sector}, whose members escape every single--tail family.
\end{proposition}
\begin{proof}
$U$ lies in the tail family of $U\setminus\{r\}$ iff $T(U)-1/r\ge2$, which is easiest for
$r=\max U$; (i) is that case ($T=2$ being impossible by rigidity, the base is admissible). In
case (ii) no removal reaches $2$, and a fortiori the two stated $59$--subsets stay below $2$.
\end{proof}

\begin{remark}[The campaign; completeness of the certificate; the arity barrier]\label{rem:campaign}
The reciprocity certificate captures \emph{every} $q$--independent congruence obstruction of a
tail family: at a prime $\ell\mid A_SB_S$ the system degenerates to the rank--one congruences
$x^2\equiv D_S$ or $y^2\equiv D_S\pmod\ell$, while at $\ell\nmid 2A_SB_S$ the conditions define
conics in $(x,q)$ or $(y,q)$ with $\sim\ell$ points, deleting at most half of the residue classes
of $q$ and never all (the $2$--adic layer pins $q\equiv5\bmod8$ without emptiness); by CRT and
Dirichlet, admissible $q$ survive any finite set of such conditions. Hence a tail family admits a
congruence kill iff $(D_S\mid\ell)=-1$ for some $\ell\mid A_SB_S$, which the Jacobi symbol
detects for $31{,}219$ bases, and partial factorisation (\texttt{code/hunt/survivor\_kill.rs}) for a
further $15{,}264$: in total $46{,}483$ of the $49{,}961$ single--tail families at level $60$ are
proven empty ($93.0\%$) and $3{,}478$ remaining open. Of the latter, $34$ have $A_S$ and $B_S$ both
prime and are emptied by Proposition~\ref{prop:immunedecide}; the other $3{,}444$ carry a balanced
${\sim}57$--digit semiprime and are factoring--hard.
The $34$ immune families are not a bottleneck; they are the tractable ones, and they are empty.
Immunity means $A_S$ is prime, and that is exactly what makes the family decidable directly.

\begin{proposition}[Immune families are decided in polynomial time]\label{prop:immunedecide}
Let $S$ be a level--$60$ base with $A_S=N_S+2D_S$ prime and $(D_S\mid A_S)=+1$. Then a solution in the
tail family $S\cup\{q\}$ forces $x^2=A_Sq+D_S$ with $x$ one of the two square roots of $D_S$ modulo
$A_S$, so the family has at most two candidate tail primes, computable by Tonelli--Shanks. All $34$
immune families at level $60$ are empty.
\end{proposition}
\begin{proof}
Write $\sigma=N_S/D_S$, so $A_S=(\sigma+2)D_S$. Proposition~\ref{K-prop:tailbound} bounds
$q\le tD_S$ with $t=\tfrac12(\sigma+\sqrt{\sigma^2-4})$, whence
$x^2=A_Sq+D_S\le\bigl((\sigma+2)t+1\bigr)D_S^2$. At level $60$ one has $\sigma>2$ and $\sigma$ close
to $2$, so $(\sigma+2)t+1<5.202$ and $x<2.281\,D_S<(\sigma+2)D_S=A_S$; over all $49{,}961$ admissible
bases the ratio $x/A_S$ is at most $0.56982$. Hence $x$ lies in $[0,A_S)$ and, $A_S$ being prime, is one
of the two roots of $X^2\equiv D_S$. Each root gives one candidate
$q=(x^2-D_S)/A_S$, to be tested for integrality, primality and the companion condition that
$B_Sq+D_S$ be a square. Carrying this out on the $34$ families, whose $A_S$ and $B_S$ are
unconditionally prime by the \textup{APRCL} and \textup{ECPP} runs recorded above, every candidate
fails \textup{(\texttt{code/immune\_decide.gp})}, either because
$q=(x^2-D_S)/A_S$ is not a prime exceeding $787$ or because $B_Sq+D_S$ is not a square; the conclusion
is therefore \textup{[P]}, not \textup{[C]}. The script re--derives the $49{,}961$ bases, the $31{,}219$
Jacobi kills and the $34$ immune families from scratch before deciding them.
\end{proof}

Primality of $A_S$ plays no role in the bound $x<A_S$, only in the ease of extracting the roots, so
the argument is general.

\begin{proposition}[The split sieve]\label{prop:splitsieve}
Let $S$ be a finite set of primes, $D=\prod S$, and let $M$ be a finite set of primes disjoint from
$S$. Suppose the pair $a=\prod(T\cup M)$, $b=\prod(S\setminus T)$ with $T\subseteq S$ is a
derivative two--cycle. Then
\[ \prod T\ \Bigm|\ \partial\bigl(\textstyle\prod(S\setminus T)\bigr),
   \qquad
   \Bigl(\textstyle\prod T\Bigr)^{2}\partial(\textstyle\prod M)
   +\partial\bigl(\textstyle\prod T\bigr)\,\partial\bigl(\textstyle\prod(S\setminus T)\bigr)=D . \]
The second identity carries no reference to the tail; the first is its reduction modulo $\prod T$,
Theorem~\ref{K-thm:structure} removing the factor $\partial(\prod T)$, and is therefore a consequence
of it rather than an additional hypothesis. Writing $\alpha=\prod T$ and $\beta=\prod(S\setminus T)$,
the divisibility is by Lemma~\ref{lem:anatomy} the system of congruences
$\sum_{p\in S\setminus T}p^{-1}\equiv0\pmod r$, one for each $r\in T$, so the expected number of
surviving $T$ is $\prod_{p\in S}(1+1/p)$ out of $2^{|S|}$. Neither statement involves $A_S$, so the
criterion applies to the families whose $A_S$ is composite, which
Proposition~\ref{prop:immunedecide} cannot reach, and to the pair sector of
Proposition~\ref{prop:sectors}(ii), where $|M|=2$ and $\partial(\prod M)=p+q$. Two consequences are
recorded. Since $\partial(\alpha)\partial(\beta)>0$ for $T$ and $S\setminus T$ both nonempty, the
identity forces $\alpha^{2}<D$. And for $X$ a set of odd primes $\partial(\prod X)\equiv|X|\pmod2$,
so $2\in T$ forces $|T|$ odd. $0$ \texttt{sorry}
\textup{(\texttt{Erdos307.split\_criterion\_gen}, \texttt{split\_snd\_gen},
\texttt{split\_snd\_of\_criterion}, \texttt{csum\_odd\_card};
\texttt{lean/Erdos307/SplitSieve.lean})}.
\end{proposition}
\begin{proof}
Leibniz on the disjoint union gives $\partial(\prod(T\cup M))=\prod M\cdot\partial(\prod T)
+\prod T\cdot\partial(\prod M)$. The equation $\partial(b)=a$ reads
$\partial(\beta)=\alpha\prod M$, which is the divisibility. The equation $\partial(a)=b$ reads
$\beta=\prod M\cdot\partial(\alpha)+\alpha\,\partial(\prod M)$; multiplying by $\alpha$, replacing
$\alpha\prod M$ by $\partial(\beta)$ and using $\alpha\beta=D$ gives the identity. For the converse
direction of the first claim, reduce the identity modulo $\alpha$: both $\alpha^{2}\partial(\prod M)$
and $D$ vanish, so $\alpha\mid\partial(\alpha)\partial(\beta)$, and
$\gcd(\alpha,\partial(\alpha))=1$ by Theorem~\ref{K-thm:structure}.
\end{proof}

\begin{remark}[What the sieve decides]\label{rem:splitsieverun}
The criterion was applied to every level--$60$ single--tail base surviving the certificate of
Proposition~\ref{prop:tailkill}, $18{,}742$ of them, a superset of the $3{,}444$ open families, over
all $T$ with $|T|\le6$: about $9\times10^{11}$ splits, no factorisation, and no two--cycle. The
observed survivor count is $5.20$ per base against the predicted $5.17$
\textup{(\texttt{code/split\_sieve.rs}, \texttt{code/gen\_bases.rs})}. The range $|T|\ge7$ is not
decided. Its expected survivor count over $7\le|T|\le34$ is about $28$, and deciding it by
enumeration costs at least $7\times10^{2}$ core--years, so the sieve reduces the level--$60$
single--tail question without closing it.
\end{remark}

\begin{proposition}[Single--tail families are decided by factoring]\label{prop:lvl60factor}
Let $S$ be a set of primes containing $2$, with $D_S=\prod S$ squarefree, and let $q$ be a prime
completing $S$ to a solution support. Then $A_S$ is odd and coprime to $D_S$, and $x=a+b$ satisfies
$x^2=A_Sq+D_S$ with $0<x<A_S$. Hence $x$ is one of the $2^{\omega(A_S)}$ square roots of $D_S$ modulo
$A_S$, each giving a single candidate $q=(x^2-D_S)/A_S$, and the family is decided by one
factorisation of $A_S$ together with $O(2^{\omega(A_S)})$ primality tests.
\end{proposition}
\begin{proof}
Since $2\in S$ and $D_S$ is squarefree, exactly one term of $N_S=\sum_{p\in S}D_S/p$ is odd, so $N_S$ is
odd and $A_S=N_S+2D_S$ is odd. If $p\in S$ divided $A_S$ it would divide $N_S$, but $N_S\equiv D_S/p
\not\equiv0\pmod p$; as the primes of $D_S$ are exactly $S$, $\gcd(A_S,D_S)=1$. For the bound, write
$\sigma=N_S/D_S$ and $A_S=(\sigma+2)D_S$. Proposition~\ref{K-prop:tailbound} gives $q\le tD_S$ with
$t=\tfrac12(\sigma+\sqrt{\sigma^2-4})$, so $x^2\le\bigl((\sigma+2)t+1\bigr)D_S^2$. Across all
$49{,}961$ admissible level--$60$ bases $\sigma\in(2,2.00236]$, whence $(\sigma+2)t+1<5.202$ and
$x<2.281\,D_S$, while $A_S=(\sigma+2)D_S>4\,D_S$; the ratio $x/A_S$ is at most $0.56982$ over the whole
set \textup{(\texttt{code/lvl60check.gp})}. Thus $x<A_S$. When instead $\sigma\le2$, as for the
pair--sector bases of Proposition~\ref{prop:sectors}(ii), Corollary~\ref{K-cor:maxbound} bounds the tail
by $q\le D_S/(2D_S-N_S)=D_S/|B_S|$, so $x^2=A_Sq+D_S\le A_SD_S/|B_S|+D_S<A_S^2$ because $|B_S|\ge1$ and
$A_S>D_S$; the conclusion $x<A_S$ therefore holds in both regimes. Being odd and coprime to $D_S$,
$A_S$ admits exactly $2^{\omega(A_S)}$ square roots of $D_S$ when $D_S$ is a residue modulo every prime
power dividing it, and none otherwise; these are computed from the factorisation by Tonelli--Shanks,
Hensel lifting at any repeated factor, and the Chinese remainder theorem.
\end{proof}

This replaces a $2^{58}$ search per family by a single factorisation, and the integer to be factored is
$A_S\approx4D_S$, of $114$ decimal digits, not the $230$--digit product the divisor search ranges over.
It also supersedes the reciprocity route rather than extending it: a congruence kill decides at best
half of the split families, since for $A=pq$ with $(D_S\mid A)=+1$ the two Legendre symbols are either
both $+1$, when no kill exists, or both $-1$, when one does, whereas the root enumeration decides every
family whose $A_S$ is factored. One cheap idea deserves to be recorded as closed, because the level--$60$ bases are far from random:
all $49{,}961$ share the $39$ forced primes below $167$, so one may hope that two of the $A_S$ share a
prime factor and split each other for free. A batch \textsc{gcd} in the sense of Bernstein (one
product tree over the hard cofactors, then a remainder tree) tests this for every pair at once in
seconds. Over the $7{,}713$ families that survive the Jacobi certificate and trial division to $10^6$,
a superset of the $3{,}478$, exactly two of the $6{,}972$ composite cofactors of $A_S$ share a factor,
and both shared factors have seven digits, so they lie barely past the trial--division bound and are
found instantly by any factoring method. The structure the bases share is in $D_S$ and does not
propagate to $A_S=N_S+2D_S$: at the wall the residual behaves as a set of pairwise coprime balanced
semiprimes, and no amount of shared structure among the bases is going to split them.
\textup{(\texttt{code/batchgcd.gp})}

Level $60$ is therefore not an exponential obstruction but a factoring
project: the $34$ with $A_S$ prime are already done, and $3{,}444$ integers of $114$ digits remain.
A $114$--digit general number field sieve is routine, costing core--days apiece rather than the
$3.2\times10^7$ core--years the divisor search of Proposition~\ref{K-prop:tailbound} would demand, a
reduction by six orders of magnitude. It should be said plainly that this does not by itself deliver
$|P\cup Q|\ge61$. The single--tail sector is only one of the two cases of
Proposition~\ref{prop:sectors}, and the pair sector contributes a further $3{,}013{,}379$ surviving
bases, each an arity--one tail family to which Proposition~\ref{prop:lvl60factor} applies equally.
Factoring those as well is a computation some three orders of magnitude larger, tens of thousands of
core--years, which is beyond what has been carried out here and is quoted as the honest price, not
as a claim of feasibility. What has changed is the character of the obstruction: level $60$ is no
longer exponential in $|S|$, but linear in the number of families, each at the cost of one
factorisation.

Those $3{,}478$ are expected to be empty, and by a wide margin: evaluating the expected count of
Remark~\ref{rem:killcount} over them, with $\kappa_S\le1$ and the tail sum bounded by
$\log\log(tD_S)$, gives at most $5.5\times10^{-108}$, the bases having $D_S$ of $113$ digits. The
divergence of the global count comes from large $a$ and not from this level, and the first solution
the heuristic locates sits at $10^{116.7}$ or beyond. So the factoring wall guards the verification of
something the heuristic already asserts, rather than a plausible solution. The weapon is intrinsically \emph{arity--one}, and provably fails on the pair sector. With
$s=p+q$, $t=pq$, a pair--sector member requires the two simultaneous squares $x^2=A_Wt+D_Ws$ and
$y^2=B_Wt+D_Ws$ (the would--be third condition $s^2-4t=\square$ is automatic, being $(p-q)^2$),
together with the reciprocity--linked constraints $(D_Wp\mid q)=(D_Wq\mid p)=1$ (mod $q$, a
solution with $q\mid a$ has $N_U\equiv D_Wp$, forcing $D_Wp$ to be a square). At a prime
$\ell\mid A_W$ the condition $x^2\equiv D_Ws\pmod\ell$ constrains only $s$, which the free second
prime absorbs; the pair system is in fact locally soluble at \emph{every} modulus (a theorem,
Proposition~\ref{prop:pairlocal} below, corroborated by a direct scan over all $\ell\le600$) so,
unlike the single--tail family, the pair sector admits \emph{no} congruence obstruction at all. It
is a genuine two--parameter surface per base, $\#307$ in miniature (Remark~\ref{rem:oneside});
reciprocity cannot touch it, and only a direct both--squares search or a descent argument can.

The sector is nevertheless an explicit finite object, and small. A level--$60$ support $U$ is covered
by the one--new--prime campaign exactly when $T(U\setminus\{\max U\})>2$, the bases of
Proposition~\ref{K-prop:close59} being the $59$--prime sets of mass exceeding $2$; the complement is
$U=R\cup\{m\}$ with $|R|=59$, $m=\max U$ and $T(R)\le2<T(R)+1/m$. The mass ladder at $k=60$ puts
every element of $R$ below $793.67\times2=1587.4$, so the complement is a search over $59$--subsets
of the $250$ primes under $1588$ with $T(R)\in(2-1/\max R,\,2]$. There are exactly
$18{,}234{,}653$ of them, the largest $\max R$ being $1583$ \textup{(\texttt{code/pairsector\_count.rs},
corroborated by an independent binned counting DP in \texttt{code/pairsector\_countdp.rs})}. Each is
an arity--\emph{one} tail family in $m$: Proposition~\ref{prop:pairlocal} denies a congruence kill to
the family with both new primes free, not to these, so the $q$--independent certificate does apply to
each of the $18.2$ million bases. The immunity of the pair sector is a consequence of its two free
parameters, and fixing one of them buys the weapon back.

That the sector is finite at all rests on $T(R)\ne2$: at $T(R)=2$ every $m$ would satisfy
$T(R)+1/m>2$ and the tail would be unbounded. It cannot occur. If $2\notin R$ then $D$ is odd and each
$D/p$ is odd, so $N=\sum_{p\mid D}D/p$ is a sum of $59$ odd terms and is odd; if $2\in R$ then $D/2$
is odd and the other $58$ cofactors are even, so $N$ is again odd. Since $2D$ is even, $N\ne2D$
\textup{(\texttt{Erdos307.mass\_ne\_two}; \texttt{code/pairsector\_parity.gp})}.

Running the certificate over the sector kills $10{,}457{,}149$ of the $18{,}234{,}653$ bases
($57.3\%$) by the Jacobi symbol alone, and trial division of $A_R$ and $B_R$ to $10^6$ closes a
further $4{,}580{,}243$ of the survivors ($58.9\%$ of them): in total $15{,}037{,}392$, or $82.5\%$,
with $3{,}197{,}261$ still open \textup{(\texttt{code/pairsector\_kill.rs}, \texttt{code/pairsector\_factor.rs}; the same
code in campaign mode reproduces the $49{,}961$, $31{,}219$ and $18{,}742$ of the paragraph above
exactly)}. A Pollard rho pass with a budget of $5\times10^4$ iterations per base then kills a further
$183{,}882$, leaving $3{,}013{,}379$ open ($83.5\%$ of the sector eliminated); the marginal return has
fallen to $5.7\%$, which is what one expects once the small factors are gone and what remains is a
balanced semiprime. The pair sector is therefore not a barrier of a different kind from the one--new--prime
families: it is the same computation at $365$ times the scale, and it responds to the same weapon.
\end{remark}

\begin{proposition}[The pair sector admits no congruence kill]\label{prop:pairlocal}
Let $W$ be any pair--sector base and $A=A_W$, $B=B_W$, $D=D_W$, so that a pair--sector solution
requires $x^2=A\,\alpha\beta+D(\alpha+\beta)$ and $y^2=B\,\alpha\beta+D(\alpha+\beta)$ with
$\alpha,\beta$ the two new primes. For every prime power $\ell^j$ this system has a solution in
units $\alpha,\beta$ and integers $x,y$ at a nonsingular point. Consequently, by CRT and Dirichlet,
no finite system of congruence conditions empties any pair--sector family: the arity barrier of
Remark~\ref{rem:campaign} is a theorem, and the certificate programme provably ends at arity one.

\textup{Formalised, in part.} \texttt{Erdos307.no\_pinning\_isUnit} is the step that makes the
tail--kill mechanism unavailable, \texttt{collapse\_identity} is regime~\textup{(2)}'s
observation that the two square conditions are one, \texttt{regime\_one\_first} and
\texttt{regime\_one\_second} are the explicit unit of regime~\textup{(1)}, and
\texttt{targets\_agree\_mod\_eight} is the mod--$8$ collapse of regime~\textup{(3)}. The lift from
$\ell$ to $\ell^j$ that regime~\textup{(1)} needs is \texttt{square\_lifts\_to\_prime\_powers},
proved from \texttt{hensel\_all\_powers}: if $\ell$ is an odd prime not dividing $c$ and $c$ is a
square modulo $\ell$, it is a square modulo every power of $\ell$. Regime~\textup{(2)} needs no citation to Weil's bound. Its two steps are
\texttt{Erdos307.pairlocal\_inner\_sum}, the evaluation of the inner sum of $S_{fg}$ as
$-\chi\bigl((A\alpha+D)(B\alpha+D)\bigr)$ by the quadratic complete sum of
Proposition~\ref{prop:localcomplete}, with its side condition discharged as $4D^2\alpha^2\ne0$; and
\texttt{pairlocal\_count\_pos}, the arithmetic turning the resulting $O(\ell)$ bounds into a positive
count from $\ell\ge17$. $0$ \texttt{sorry} \textup{(\texttt{lean/Erdos307/PairLocal.lean},
\texttt{lean/Erdos307/PairLocalCount.lean})}. Checked directly: over $872$ random bases with
$A-B=4D$ and every odd $\ell\le400$, the count of admissible $(\alpha,\beta)$ was positive in every
case and tracks $(\ell-1)^2/4$; the crude bounds $|S_f|,|S_g|\le2\ell$, $|S_{fg}|\le3\ell$ held in
all $409$ cases tested \textup{(\texttt{code/pairlocal\_count.gp})}.
\end{proposition}
\begin{proof}
First, no $\ell$ pins either target to a constant: that needs $\ell\mid A$ (or $B$) \emph{and}
$\ell\mid D$, whence $\ell\mid N$, contradicting $\gcd(N,D)=1$ (rigidity); so the tail--kill
mechanism is unavailable, and solubility is decided in three regimes.

\emph{(1) $\ell$ odd, $\ell\mid D$.} This covers every odd $\ell\le103$: at level $60$, omitting a
prime $p$ caps $T(U)$ at $T_{61}-1/p$, and $1/p>T_{61}-2=0.00944\ldots$ exactly for $p\le103$, so
all such primes lie in $U$, hence in $W$ (the two adjoined primes are the largest). Take $\beta=1$,
$x=1$, $\alpha=(1-D)/(A+D)\bmod\ell^j$, a unit since $A+D\equiv N\not\equiv0\pmod\ell$: the first
equation holds by construction, and the second target $(B+D)\alpha+D\equiv N\cdot N^{-1}=1\pmod\ell$
is a unit congruent to a square, so $y$ exists modulo $\ell^j$.

\emph{(2) $\ell$ odd, $\ell\nmid2D$} (in particular every $\ell\ge107$ outside $W$). The two
conditions are secretly \emph{one}: setting $m=(x^2-y^2)/4D$ and $u=(x^2-Am)/D$, the relation
$A-B=4D$ gives the polynomial identity $Bm+Du=y^2$. Only \emph{existence} is wanted, and counting the
pairs $(\alpha,\beta)$ directly settles it with no curve and no Weil bound. Write $\chi$ for the
quadratic character, $f=A\alpha\beta+D(\alpha+\beta)$ and $g=B\alpha\beta+D(\alpha+\beta)$. Summing
$\tfrac14(1+\chi(f))(1+\chi(g))$ over $\alpha,\beta\in\mathbb F_\ell^{\times}$ and discarding the at
most $2(\ell-1)$ pairs with $fg=0$,
\[ 4\,\#\{\text{both nonzero squares}\}\;\ge\;(\ell-1)^2+S_f+S_g+S_{fg}-8(\ell-1). \]
For fixed $\alpha$ each target is \emph{linear} in $\beta$, so the inner sums of $S_f$ and $S_g$
vanish by the linear evaluation above, leaving $|S_f|,|S_g|\le2\ell$. And $fg$ is a \emph{product of
two linear forms} in $\beta$, so its inner sum is the quadratic evaluation, whose side condition
$ad\ne bc$ is here $D\alpha(A-B)\alpha=4D^2\alpha^2\ne0$, exactly the hypotheses
$\ell\nmid2D$ and $\alpha\ne0$; summing the resulting $-\chi\bigl((A\alpha+D)(B\alpha+D)\bigr)$ over
$\alpha$ is that same evaluation again, so $|S_{fg}|\le3\ell$. Hence the count is positive as soon as
$(\ell-1)^2-7\ell-8(\ell-1)>0$, i.e.\ from $\ell\ge17$, far below the $\ell\ge107$ in play. The point
is nonsingular ($\alpha\beta\ne0$; the Jacobian contains $\operatorname{diag}(2x,2y)$), so it lifts
to every $\ell^j$.

\emph{(3) $\ell=2$.} Both targets are odd and agree modulo $8$ (as $\alpha\beta$ is odd,
$D(\alpha+\beta)\equiv0\bmod4$, and $(A-B)\alpha\beta=4D\alpha\beta\equiv0\bmod8$), and a finite
table over $(A\bmod8,\,D\bmod8)$ shows that for every base some achievable pair
$(\alpha\beta,\alpha+\beta)\bmod8$ makes the common target $\equiv1\pmod8$; a $2$--adic square,
soluble modulo $2^j$ for every $j$.

All three regimes, the collapse identity, the threshold $1/103>T_{61}-2>1/107$, and the mod--$32$
layer were verified computationally on three bases and every $\ell\le600$
(\texttt{code/pairlocal.py}; the measured density of admissible $(x,y)$ in regime~(2) is
$0.497$--$0.503$ against the predicted $\tfrac12$).
\end{proof}

\begin{remark}[Anatomy of the certificate; the empirical $5/8$ law]\label{rem:fiveeighths}
Writing $D=2D_1$, the reciprocity signs in the certificate cancel
\textup{(}$\tfrac{A-1}2+\tfrac{N-1}2=N-1+2D_1$ is even\textup{)}, leaving the clean product
\[ \left(\tfrac{D}{A}\right)\;=\;\left(\tfrac{2}{A}\right)\left(\tfrac{D_1}{N}\right),\qquad
\left(\tfrac{2}{A}\right)=+1\iff N\equiv3,5\ (\mathrm{mod}\ 8), \]
and one layer down $(N\mid D_1)=\prod_{p\mid D_1}\bigl(\tfrac{D/p}{p}\bigr)$, since
$N\equiv D/p\pmod p$: the cofactor structure of rigidity recurs inside its own certificate. This product evaluates in closed form: with $D/p=2D_1/p$ and quadratic
reciprocity,
\[ (N\mid D_1)=(2\mid D_1)\!\!\prod_{\{p,q\}\subset D_1}\!\!(-1)^{\frac{p-1}2\frac{q-1}2}
=(-1)^{\,s+\binom{t}{2}},\quad s=\#\{p\mid D_1:p\equiv3,5\ (8)\},\ \ t=\#\{p\mid D_1:p\equiv3\ (4)\}, \]
a R\'edei genus character of $D_1$, \emph{independent of $N$}. This is forced, not merely observed:
$(N\mid D_1)=\prod_{p\mid D_1}(N\mid p)$, and cofactor survival gives $N\equiv D/p=2D_1/p\pmod p$, so
$(N\mid p)=(2\mid p)\,(D_1/p\mid p)$. The first factors multiply to $(-1)^{s}$ since $(2\mid p)=-1$
exactly for $p\equiv3,5\pmod 8$; and, each unordered pair $\{p,q\}$ occurring once in each of the two
cofactors,
\[ \prod_{p\mid D_1}\Bigl(\tfrac{D_1/p}{p}\Bigr)=\prod_{\{p,q\}\subset D_1}\Bigl(\tfrac qp\Bigr)
\Bigl(\tfrac pq\Bigr)=\prod_{\{p,q\}\subset D_1}(-1)^{\frac{p-1}2\frac{q-1}2}=(-1)^{\binom t2}, \]
by quadratic reciprocity, since $\tfrac{p-1}2$ is odd exactly for $p\equiv3\pmod4$. (Checked
independently, prime by prime, on $20{,}000$ bases: $0$ violations.) The inner symbol is thus not a coin but a genus invariant; its empirical fairness
($49.9\%$, $50.0\%$ over the $49{,}961$ bases) is that character's equidistribution; and
$(D_1\mid N)$ follows by the reciprocity sign $(-1)^{\frac{D_1-1}2\frac{N-1}2}$, hence depends on $N$
only through $N\bmod4$. In fact the entire certificate closes into a finite functional. Let
$v(S)=(n_1,n_3,n_5,n_7)\bmod4$ count the odd primes of $S$ in the residue classes $1,3,5,7$ modulo
$8$. Every factor above is a function of $v$ alone: $s\equiv n_3+n_5$, $t=n_3+n_7$ (with
$\binom t2$ needing only $t\bmod4$), $D_1\equiv3^{n_3}5^{n_5}7^{n_7}\pmod8$,
$S(D_1)\equiv(n_1+n_5)-(n_3+n_7)\pmod4$, and (by the mod--$8$ law of
Proposition~\ref{C-prop:mod8}, $D_1'\equiv D_1S(D_1)\pmod8$) also
$N=D_1+2D_1'\equiv D_1+2\bigl(D_1S(D_1)\bmod4\bigr)\pmod8$. Hence
\[ (D\mid A)\;=\;K\bigl(v(S)\bigr): \]
\emph{the kill is decided by sixteen residue counts modulo $4$.} The realized vectors are exactly
the coset $n_1+n_3+n_5+n_7\equiv58\equiv2\pmod4$ (all $64$ of its classes occur), $K$ is constant
on each, and the functional reproduces the direct $113$--digit Jacobi computation on every one of
the $49{,}961$ bases (\texttt{code/kill\_classvector.py}). The five--eighths law is then an exact
finite statement, and its mechanism is the \emph{reciprocity switch}: for $N\equiv1\pmod4$ the
reciprocity sign is inert and the kill reduces to the genus coin; exact conditional rate
$12{,}699/24{,}975=0.5085$; while for $N\equiv3\pmod4$ the sign couples the parity of $t$ into
the product and the rate jumps to $18{,}520/24{,}986=0.7412$ (by class modulo $8$:
$0.509,\,0.741,\,0.508,\,0.741$ at $N\equiv1,3,5,7$; the $3/4$ classes are $N\equiv3\pmod4$, not
the $(2\mid A)$ classes). The blend $\tfrac12\cdot\tfrac{24975}{49961}+\tfrac34\cdot\tfrac{24986}{49961}
=0.62503$ against the exact $31{,}219/49{,}961=0.62487$ is the whole of the ``coincidence''. The
correlation is an ensemble property, a random--subset proxy gives ${\approx}0.45/0.55$; and
the functional yields no second certificate: it recomputes the same symbol.
\end{remark}

\begin{remark}[Why a one--sided mass argument cannot reach the barrier]\label{rem:twosided}
The threshold $\sigma(U)=t+1/t$ of Proposition~\ref{prop:barthreshold} makes the barrier equivalent to a
lower bound $|\sigma(a)-1|>c$ with $c$ absolute, and it is worth recording what such a bound cannot
be proved from. The primary pseudoperfect numbers satisfy the one--sided equation $n'=n-1$ exactly,
so along them $|\sigma(n)-1|=1/n$ tends to $0$: the eight known members give
$1/n$ from $\tfrac12$ down to $1.18\times10^{-31}$. They are solutions of the coprime relaxation
(Corollary~\ref{K-cor:threethirteen}, the case $u=1$). Hence no argument that uses only $a'=b$ with $b$
squarefree and coprime to $a$ can bound $|\sigma(a)-1|$ below by an absolute constant, since the same
argument would prove the corresponding statement for the relaxation, which is false. Any such bound must
consume the second equation $b'=a$. This is why the mass arguments of Section~\ref{K-sec:barrier} of the core note, and the
Sylvester--type extremal bounds, reproduce $60$ and stop: they are one--sided, and $1/a$ is the most a
one--sided argument can give.
\end{remark}

\begin{remark}[The kill--program is one--sided]\label{rem:oneside}
Every kill is a proof that no solution hides in that family, but if \#307 has a solution, the
solution's own family is unkillable, so no accumulation of kills can overshoot the truth: the
program's asymptote \emph{is} the existence question. Together with the self--similarity of
Remark~\ref{rem:tier60} (each rung past $60$ contains the whole problem), this locates the
campaign precisely: it corners where a solution could live, one certified region at a time, and
can terminate only if \#307 has answer \textsc{no}, which the heuristics disfavour.
\end{remark}

\begin{remark}[The reciprocity kills are the local factors of the expected count]\label{rem:killcount}
The divergent naive count of Section~\ref{sec:anatomy} and the reciprocity sieve of
Proposition~\ref{prop:tailkill} are the \emph{same} computation. Decompose the expected number of
solutions over tail families $S\cup\{q\}$: the summand for a base $S$ is
$\kappa_S\,\bigl(D_S\sqrt{\mathrm{mass}^2-4}\bigr)^{-1}\sum_q q^{-1}$, where
$\kappa_S=\prod_\ell(\text{local solubility density at }\ell)$ is exactly the reciprocity data.
A reciprocity--killed base carries an inert $\ell\mid A_S$ with $(D_S\mid\ell)=-1$, so
$N_S+2D_S\equiv D_S\pmod\ell$ is a non--residue for \emph{every} $q$, forcing $\kappa_S=0$; the
family drops from the count entirely. An immune base ($A_S,B_S$ both prime, both residues) has no
inert prime, so $\kappa_S>0$ and, treating $q$ as ranging freely with density $1/q$, the summand
\emph{looks} like a divergent $\sum_q q^{-1}$ (verified on an explicit $114$--digit immune base: local
density $1$ at every tested prime, partial sums growing without bound). That model overstates the
space: Proposition~\ref{K-prop:tailbound} bounds $q$ itself, so the true sum is finite, and
Proposition~\ref{prop:lvl60factor} identifies it exactly with the roots of $D_S$ modulo $A_S$, of which
there are only two once $A_S$ is prime. The immune part of the divergence was therefore never genuine,
only an artifact of summing a local density over an unbounded range that
Proposition~\ref{K-prop:tailbound} had already closed; the kills remove ${\sim}93\%$ of the count's mass,
and Proposition~\ref{prop:immunedecide} removes the rest of the $34$ conditionally immune families'
contribution unconditionally, both candidate values of $q$ failing outright. What
Remark~\ref{rem:killcount} originally located as the entire ``weakly favours \textsc{yes}'' signal at
level $60$ is now zero there. What remains is the pair sector, and there the picture does not simplify
the same way: the object Proposition~\ref{prop:pairlocal} proves permanently immune to any finite
congruence condition is the one with \emph{both} new tail primes free, which the campaign never
enumerates as such, fixing one of the two, as Remark~\ref{rem:campaign} does to get a finite, countable
search, buys the congruence weapon back for the resulting $18{,}234{,}653$ bases, and their $83.5\%$
elimination is a plain factoring--budget shortfall, the same story as the single--tail sector's, not an
instance of Proposition~\ref{prop:pairlocal}'s theorem. Proposition~\ref{prop:lvl60factor} applies to
that computed population exactly as it does to the single--tail residual (its proof addresses the
$\sigma\le2$ case via Corollary~\ref{K-cor:maxbound}), so the residual open question there is, again, the
single task of factoring a $\sim\!114$--digit integer per family. Proposition~\ref{prop:pairlocal}'s
genuinely uncomputed object is a separate matter: it is why the campaign stops at fixing one tail prime
rather than a limitation discovered by trying and failing to go further. This does not touch the
reliability of the count (which ignores the $N(r)\le1$ rigidity), but it retires the level--$60$
half of what it once pinned down as the heuristic's \textsc{yes}.
\end{remark}

\begin{proposition}[Local completeness of the reciprocity sieve]\label{prop:localcomplete}
For an admissible base $S$ and an odd prime $\ell$, the admissible tail residues
$Q_\ell=\{q\bmod\ell:(A_Sq+D_S\mid\ell)\ne-1\text{ and }(B_Sq+D_S\mid\ell)\ne-1\}$ fall into exactly three
regimes. \textup{(i)} $\ell\mid D_S$, which covers every $\ell\le167$ uniformly over the $49{,}961$
admissible bases, since $1/p\ge(T_{59}-2)+\tfrac1{281}$ forces every prime $p\le167$ into every base:
then $A_S\equiv B_S\equiv N_S\not\equiv0\pmod\ell$, the two conditions collapse to the single condition
$(N_Sq\mid\ell)\ne-1$, and $|Q_\ell|=(\ell+1)/2\ni0$. \textup{(ii)} $\ell\mid A_SB_S$: one condition is
the constant $(D_S\mid\ell)$, so $Q_\ell=\emptyset$ iff $(D_S\mid\ell)=-1$ (exactly the certificate of
Proposition~\ref{prop:tailkill}) and otherwise $|Q_\ell|\ge(\ell-1)/2$. \textup{(iii)}
$\ell\nmid A_SB_SD_S$ (hence $\ell\ge173$): the complete character sums give
$|Q_\ell|\ge(\ell-3)/4-2\ge40$ (the linear sums vanish, the quadratic sum is $-\chi(A_SB_S)$, boundary
terms $\le2$). \textup{(}Regimes~\textup{(i)} and~\textup{(ii)} are formalised:
\texttt{Erdos307.regime\_one\_collapse}, \texttt{regime\_one\_same\_argument} and
\texttt{regime\_two\_constant} \textup{(}\texttt{lean/Erdos307/LocalComplete.lean}\textup{)}.
Regime~\textup{(iii)}'s counting step is \texttt{regime\_three\_count}, and its two complete
character--sum hypotheses are now theorems rather than citations:
$\sum_q\chi(aq+b)=0$ for $a\ne0$ is \texttt{Erdos307.sum\_quadraticChar\_affine}, the reindexing
$q\mapsto aq+b$ followed by \texttt{quadraticChar\_sum\_zero}; and
$\sum_q\chi((aq+b)(cq+d))=-\chi(ac)$ for $ad\ne bc$ is
\texttt{Erdos307.sum\_quadraticChar\_quadratic}, proved by factoring out $\chi(ac)$, translating to
$\sum_u\chi\bigl(u(u+e)\bigr)$ with $e\ne0$, writing $u(u+e)=u^2(1+e/u)$ for $u\ne0$ so that the
character value is $\chi(1+e/u)$, and reindexing by the involution $u\mapsto e/u$ of the nonzero
elements to reach $\sum_{t\ne0}\chi(1+t)=0-\chi(1)=-1$. Both hold over any finite field of odd
characteristic. $0$ \texttt{sorry}
\textup{(\texttt{lean/Erdos307/CompleteSums.lean})}. Checked against direct evaluation for all odd
primes $\ell\le59$: the linear sum over every $(a,b)$ with $a\ne0$, and the quadratic sum over all
$33{,}008{,}952$ quadruples with $ad\ne bc$, $0$ violations
\textup{(\texttt{code/complete\_sums\_check.gp})}.\textup{)}
Consequently $Q_\ell\ne\emptyset$ forces $|Q_\ell|\ge2$ at every modulus, the surviving
local density is positive (e.g.\ $\prod_{\ell\le300}|Q_\ell|/\ell\approx2.2\times10^{-19}>0$ on the
canonical base), and by CRT \emph{no finite system of congruences (no covering) can annihilate an
unkilled family}: the binary certificate of Proposition~\ref{prop:tailkill} is the complete local theory
of the tail sector, there is no ``second kill layer,'' and the one--sidedness of
Remark~\ref{rem:oneside} is structural rather than an artefact of method. (All three regimes verified
exhaustively on the canonical base: collapse and $|Q_\ell|=(\ell+1)/2$ at
$\ell\in\{3,5,7,101,277\}$; the Weil floor at every prime $281\le\ell\le600$; $A_S,B_S$ carry no factor
below $10^6$, consistent with the balanced--semiprime wall of Remark~\ref{rem:campaign}.)
\end{proposition}

\begin{proposition}[What a barrier improvement must supply]\label{prop:barthreshold}
Let $(a,b)$ be a solution with $a<b$ and put $t=b/a$, so that $t=\sigma(a)$ and, by
$\sigma(a)\sigma(b)=1$,
\[ \sigma(U)\;=\;\sigma(a)+\sigma(b)\;=\;t+\tfrac1t. \]
For $n$ with $T_n>2$ let $t_n>1$ be the root of $t+1/t=T_n$. Since $\sigma(U)\le T_{|U|}$,
\[ \sigma(a)>t_n\ \Longrightarrow\ |U|\ \ge\ n+1, \]
and this is the only route by which a mass argument can raise the barrier. Numerically
\[
\begin{array}{l|cccccc}
n & 59 & 60 & 61 & 62 & 63 & 64\\\hline
T_n & 2.0023502 & 2.0059089 & 2.0094424 & 2.0128554 & 2.0161127 & 2.0193282\\
t_n & 1.0496677 & 1.0798804 & 1.1020081 & 1.1199914 & 1.1352477 & 1.1490254
\end{array}
\]
\end{proposition}

So the barrier level is decided entirely by how close $\sigma(a)$ may lie to $1$, and each further
level demands a definite gap: $|U|\ge61$ needs $\sigma(a)>1.0799$, a full $8\%$ above $1$.

That is why the mass mechanism is exhausted here rather than merely difficult. The only
unconditional lower bound on the gap is $\sigma(a)-1=(a'-a)/a\ge1/a$, which is worthless at
$a>10^{56}$; and $\sigma$ genuinely comes close to $1$ from above on very few primes, the
Sylvester-type set $\{2,3,7,41\}$ giving $\sigma=1.00058\ldots$ with four. Indeed no bound in
$\omega(a)$ alone can exist: the greedy Sylvester construction gives, for each $k$, squarefree $a$
with $\omega(a)=k$ and $0<|\sigma(a)-1|$ doubly exponentially small in $k$, from below
$1-\sigma=7.8\times10^{-592}$ at $k=12$, and from above $\sigma-1=4.6\times10^{-591}$ at the same
$k$ \textup{(\texttt{code/sigma\_gap.gp})}. So $|\sigma(a)-1|\ge f(\omega(a))$ is false for every
positive $f$, and the unconditional $1/a$ is of the right shape. Nothing in the cycle
equations forbids $a'-a=1$. A barrier past $60$ must therefore come from a mechanism that bounds
$\sigma(a)$ away from $1$ without appealing to the size of $U$, which is what
Theorem~\ref{K-thm:barrier} does in the opposite direction, and we know of none.

The same reduction locates the extreme case. Writing $c=b-a$, the identity
$\sigma(U)-2=(t-1)^2/t=c^2/D$ is the minus layer $D'-2D=c^2$, and $\gcd(c,D)=\gcd(b-a,ab)=1$ because
$\gcd(a,b)=1$. Since an admissible support contains every prime $\le167$
\textup{(}Remark~\ref{rem:tier60}\textup{)}, $c$ is divisible by none of them, so $c=1$ or
$c\ge173$. The branch $c=1$ is $D'=2D+1$, equivalently $a'=a+1$ and $(a+1)'=a$; it is the closest a
solution could come to $\sigma(a)=1$, and it is exactly the configuration a mass argument cannot
reach. Below $3\times10^{6}$ the equation $a'=a+1$ has only the solutions $30$, $858$, $1722$ and
$66{,}198$, and none extends: for the first three $a+1$ is prime, so $(a+1)'=1$, and for the last
$a+1$ is not squarefree. The branch is not excluded, only unsearched, since the barrier puts it above
$10^{56}$. The level--$59$
computation that gives $|U|\ge60$ is, in these terms, an exhaustive verification that
$\sigma(a)>t_{59}$, and Proposition~\ref{K-prop:levelbarrier} says the same verification is unavailable
one level up.

\begin{remark}[The tail is not a Lehmer question]\label{rem:notlehmer}
Remark~\ref{B-rem:whatisleft} names the exponential Pell orbit of Remark~\ref{rem:lehmer} as the only live branch, and it is natural to read the surviving question as the primality of a term $q_n$ along it. That reading does not hold. The orbit is genuinely infinite and both square conditions do hold along all of it; what fails is that no term past $N$ can be prime, by Proposition~\ref{K-prop:tailbound}. The single identity responsible is $A-B=4D$, used elsewhere only to form the Pell
system and never to factor $4Dq$.
\begin{remark}[The threshold in largest--prime coordinates, and which sectors it leaves alive]\label{rem:sectorlive}
Remark~\ref{rem:tier60} already writes a cycle as $a=qm$ with $q$ the largest prime; Bado
\cite{Bado2026} makes the same reduction exact and independent of the level. With $q=P^{+}(ab)$,
$q\mid a$, write $a=qd$, $b=e$: then $e=d+qd'$, $e'=qd$, and $q=(e-d)/d'$. Here
$\sigma(a)=1/q+\sigma(d)$, so $t$ (or $1/t$, according to which side carries $q$; $t+1/t$ is
symmetric and the reading below is unaffected) is $\sigma(d)$ up to $1/q$, and
Proposition~\ref{prop:barthreshold} becomes a statement about the \emph{sector} $d$. It is very
uneven. \textup{(}The rows below take $q$ astronomically large, which
Proposition~\ref{K-prop:close59} supplies once $d$ is small: at $q=281$ the first row would read
$1295$.\textup{)}

\smallskip
\begin{center}
\begin{tabular}{lccc}
$d$ & $\sigma(d)$ & $\sigma+1/\sigma$ & $|U|\ge$\\ \hline
$2$ & $0.500000$ & $2.500000$ & $1413$\\
$2\cdot3$ & $0.833333$ & $2.033333$ & $69$\\
$2\cdot3\cdot7$ & $0.976190$ & $2.000581$ & $59$\\
$2\cdot3\cdot7\cdot43$ & $0.999446$ & $2.000000$ & $59$\\
$2\cdot3\cdot11\cdot23\cdot31$ & $0.999979$ & $2.000000$ & $59$\\
$3$ & $0.333333$ & $3.333333$ & ${\sim}1.1\times10^{8}$\\
\end{tabular}
\end{center}
\smallskip

The last row is a Mertens estimate, the exact $K$ being past $\pi(10^6)$. Inverting
$t+1/t\le T_K$ in the same convention gives the windows: $|U|\le60$ forces
$\sigma(d)\in[0.9260,1.0799]$, the $t_{60}$ of the table above, and $|U|\le68$ forces
$\sigma(d)\in[0.8372,1.1945]$ \textup{(\texttt{code/sector\_window.gp})}. So a sector can be cheap
only if $\sigma(d)$ sits within a few per cent of $1$. The primary pseudoperfect numbers, having
$\sigma(n)=1-1/n$ exactly, are among those extremal sectors (they are not the only ones, the
Sylvester--type $\{2,3,7,41\}$ above being another) and by
Corollary~\ref{K-cor:threethirteen} they are the $u=1$ solutions of \#307's coprime relaxation, so
Erd\H{o}s~\#313 sits at the floor of this barrier as well.

The reduction is not new to the tail theory: it \emph{is} the parametrisation of
Proposition~\ref{K-prop:tailbound}, with $\alpha=d$ and $\beta=e$. Substituting $e=d+qd'$ and $e'=qd$
into $\alpha^2q^2-Nq+(\beta^2-D)$ at $\alpha=d$ gives $e(e-d-qd')=0$ identically
\textup{(\texttt{code/alpha\_is\_d.gp}; the discriminant identity is $4d^2$ times the same
vanishing, not a second fact)}. Two things follow. Everything proved about tail families is
a statement about sectors: the reciprocity kill of Proposition~\ref{prop:tailkill}, and the
factoring wall behind the families it leaves open. And the sector $d=2$ is $\alpha=2$, which the
table puts at $|U|\ge1413$, so the level--$60$ programme never meets it. Negatively, the reduction
does not shorten that residual: the unknown is the divisor $\alpha\mid D$, of which there are
$2^{59}$, and the only $\alpha$--independent certificate is the Jacobi one already used.
\end{remark}

\begin{remark}[The sector barrier sharpened by $d'$, and a second derivation of $60$]\label{rem:sectordprime}
The sector reading above uses only the mass. It can be sharpened, and the sharpening is what
Bado \cite{Bado2026b} exploits at $d=42$. Since $e=d+qd'$, one has $\gcd(e,d)=\gcd(qd',d)=1$ and
$\gcd(e,d')=\gcd(d,d')=1$, so
\[ \operatorname{supp}(e)\ \text{avoids}\ \operatorname{supp}(d)\cup\operatorname{supp}(d')
   \qquad\text{in every sector,} \]
and the mass $\sigma(e)$, which is $1/\sigma(d)$ to many places, must be carried by primes drawn
from a thinned set. Two further identities come with it, both immediate from $e'=qd$:
$de-d'e'=d^2$ and $\sum_{r\mid e}1/r=d/d'-d^2/(d'e)$; and the associated defect
$\Delta_d(R)=dR-d'R'$ satisfies $\Delta_d(Rr)=r\Delta_d(R)-d'R$ with terminal value
$\Delta_d(e)=d^2$ \textup{(\texttt{code/sector\_general.gp}, symbolic in $d,d',q,R,R'$)}. When $d$
is primary pseudoperfect, $d'=d-1$; that is the only reason those sectors look distinguished.

Running the resulting bound sector by sector (mass $d/d'$ from primes avoiding
$\operatorname{supp}(d)\cup\operatorname{supp}(d')$, together with the parity constraint that $d$
even forces $e$ odd and $\omega(e)$ even) gives
\[ \min_{d}\ \bigl(1+\omega(d)+\omega(e)\bigr)\ =\ 60, \]
over all sectors $d$ built from primes up to $47$, attained at $d=15470=2\cdot5\cdot7\cdot13\cdot17$
and at many others, among them the primary pseudoperfect $47058=2\cdot3\cdot11\cdot23\cdot31$
\textup{(\texttt{code/sector\_dprime.gp})}. This is Theorem~\ref{K-thm:barrier}'s constant recovered
by a route sharing nothing with its proof: that one is an exhaustive computation over the $49{,}961$
admissible $59$--prime supports, this one a largest--prime descent and a reciprocal--mass count. It
does not improve the barrier, it bottoms out at exactly $60$, but two independent derivations
of the same constant are worth recording, and the sectors where it is tight are where any
improvement must act.

At $d=42$ the mass bound alone gives $\omega(e)\ge62$ and $|U|\ge66$; \cite{Bado2026b} obtains
$\omega(e)\ge64$ and $|U|\ge68$ by adjoining the residue conditions
$\sum_{r\mid e}r^{-1}\equiv0\pmod 3$ and $\pmod 7$ and $\prod_{r\mid e}r\equiv1\pmod{41}$ and
maximising over admissible supports. We have re-run that certificate independently, obtaining the
same maximum and the same $51{,}281$ reachable states. The residue sieve cannot be pushed further as
a sieve: the complete system of local conditions is the sector equation itself. For each $r\mid e$,
reducing $e'=qd$ modulo $r$ gives $\prod_{s\mid e,\,s\ne r}s\equiv-d^2/d'\pmod r$, one condition
per prime of $e$; conversely these, with $e\equiv d\pmod{d'}$ and $1<\sigma(e)<2$, force $e'=qd$
exactly, both sides lying in $(e,2e)$. Bado's conditions modulo $3$ and $7$ are the sums of these;
the maximiser certifying his bound violates them at five of its six smallest primes, and adjoining
the three at $5,11,13$ to the dynamic programme lowers the $64$--prime maximum from $1.03032$ to
$1.02853$, still above $42/41$ \textup{(\texttt{code/sector42\_znam\_violations.gp},
\texttt{code/sector42\_znam\_dp.rs})}. The next step is therefore not a sieve but an exact
enumeration, and it goes through.
\end{remark}

\begin{proposition}[The $d=42$ sector at $64$ primes]\label{prop:sector42}
No solution of \#307 has $a=42q$ with $q$ prime and $\omega(b)=64$. Hence in that sector
$\omega(b)\ge66$ and $|P\cup Q|\ge70$.
\end{proposition}

\begin{proof}[Proof (finite verification)]
Write $e=b$, $T=42/41$, and let $A_1<A_2<\cdots$ be the primes other than $2,3,7,41$, $S_j=\sum_{i\le j}1/A_i$.
By the sector identity, $\sigma(e)=T-1764/(41e)$, and $e>3.5\times10^{57}$ by Theorem~\ref{K-thm:barrier},
so $T-10^{-55}<\sigma(e)<T$. Exactly, $A_{64}=337$, $A_{197}=1229$, $A_{198}=1231$, and
$S_{61}+3/1231<T$ \textup{(\texttt{code/sector42\_k64.gp})}. If $m$ of the $64$ primes of $e$ exceed
$1229$ then $\sigma(e)\le S_{64-m}+m/1231$, which for $m\ge3$ is below $S_{61}+3/1231<T-10^{-55}$; so
$m\le2$. Next, if $\operatorname{supp}(e)=S\cup\{\ell\}$ with $R=\prod S$, the identity
$42e-41e'=1764$ reads $\ell\,\Delta_{42}(R)=41R+1764$, so $\Delta_{42}(R)>0$ and $\ell$ is the integer
$(41R+1764)/\Delta_{42}(R)$: the last prime is determined by the others.

\emph{Case $m\le1$.} Take $\ell$ the prime exceeding $1229$ if $m=1$, else the largest prime of $e$, so
$\ell\ge A_{64}$. The other $63$ primes $S$ lie among $A_1,\dots,A_{197}$ with
$\sigma(S)=\sigma(e)-1/\ell\in[T-1/A_{64}-10^{-55},T)$. The enumeration tests $36{,}901{,}596{,}763$ sets, every
such $S$ among them; for none is $(41R+1764)/\Delta_{42}(R)$ an integer.

\emph{Case $m=2$.} Let $\ell_2\le\ell_1$ be the two primes exceeding $1229$, $S$ the other $62$, and
$\delta=T-\sigma(S)$. Then $1/\ell_1+1/\ell_2\in(\delta-10^{-55},\delta)$, so $\delta<2/1231$ and
$\ell_2\in(1/\delta,\,2/\delta+1)$. The enumeration tests $166{,}213$ sets, every such $S$ among them; the least $\delta$ among them is
$1.539\times10^{-6}$, so every $\ell_2$ range is finite (largest upper bound $1{,}299{,}387$). For each
$S$ and each prime $\ell_2$ in its range, $S\cup\{\ell_2\}$ determines $\ell_1$ as above; it is never
an integer.

The enumeration prunes on the mass in floating point with a margin of $10^{-9}$ on both ends of each
window, so no set in a window is lost; the integrality test is a double--double prefilter with a
conservative band followed by exact integer arithmetic ($7{,}454{,}365$ and $2{,}495$ exact tests
respectively), the integer layer checked against PARI to the digit
\textup{(\texttt{code/sector42\_k64.rs}, two runs with independent big--integer implementations)}. The
non--enumerative half of the argument (the forced--prime identity and its converse, the sign
condition, the case--split inequality and the parity step) is formalised in
\texttt{lean/Erdos307/Sector42.lean}, with the two finite computations appearing as explicit
hypotheses of the concluding statement.
So $\omega(e)\ne64$. By \cite{Bado2026b}, $\omega(e)\ge64$ and $\omega(e)$ is even ($e$ is odd and
$e'=42q$ is even); hence $\omega(e)\ge66$. Finally $q\notin\{2,3,7\}$ and $q\nmid e$, since $\gcd(a,b)=1$ and
$a=42q$, so the four listed primes are distinct and $|P\cup Q|=|\{2,3,7,q\}|+\omega(e)\ge70$.
\end{proof}

\begin{proposition}[The primary pseudoperfect sectors]\label{prop:ppnsectors}
Let $d$ be squarefree and $\operatorname{supp}(e)=S\cup\{\ell\}$ with $R=\prod S$. The sector
equation $de-d'e'=d^2$ reads
\[ \ell\bigl(dR-d'R'\bigr)\;=\;d'R+d^2, \]
so the last prime is determined by the others in every sector, not only at $d=42$. Running the
resulting enumeration at the first admissible $\omega(e)$ above the mass floor gives, in each case,
no solution:
\[\begin{array}{lrrrl}
d & \omega(e) & \text{sets} & |P\cup Q|\ \text{before} & \text{after}\\[2pt]
42 & 62 & 175 & 66 & 68\\
47058 & 54 & 41{,}654 & 60 & 62\\
2214502422 & 70 & 11{,}646 & 77 & 79
\end{array}\]
\textup{(\texttt{code/sector\_kexclude.rs}; the identity is checked symbolically and in exact
rationals in \texttt{code/sector\_kexclude.gp})}. At each $\omega(e)$ in the table the mass bound
$S_{\omega-1}+1/A_{N+1}<d/d'$ excludes even one prime beyond the truncation, so there a single phase
is a complete argument. At $d=47058$ this raises a primary pseudoperfect sector off the global
barrier: that sector sat at exactly $60$, and $|P\cup Q|\ge62$ there. With
Proposition~\ref{prop:sector42} the $d=42$ sector stands at $|P\cup Q|\ge70$.

The condition that makes a phase count sufficient is worth isolating, because it prices every rung of
the ladder in advance and because a phase--$1$ enumeration that returns nothing proves nothing without
it.

\begin{proposition}[When the sector enumeration is complete]\label{prop:phasecount}
Let $d$ be squarefree, $d'=\sum_{p\mid d}d/p$, $t=d/d'$, and let $A_1<A_2<\cdots$ be the primes
dividing neither $d$ nor $d'$, with $S_m=\sum_{i\le m}1/A_i$. Fix a truncation $N$ and a target
$\omega(e)=k$. If $j$ of the primes of $e$ exceed $A_N$ then
\[ \sigma(e)\;\le\;S_{k-j}+\frac{j}{A_{N+1}}, \]
so whenever $S_{k-j}+j/A_{N+1}<t$ no admissible $e$ has exactly $j$ primes beyond the truncation.
Writing $J=J(k,N)$ for the least such $j$, every admissible $e$ has at most $J-1$ primes beyond
$A_N$, and the phases $1,\dots,J-1$ exhaust the possibilities: the enumeration is complete.
\end{proposition}
\begin{proof}
The $k-j$ primes of $e$ not exceeding $A_N$ lie among $A_1,\dots,A_N$, so their reciprocals sum to at
most that of the $k-j$ smallest, namely $S_{k-j}$; each of the remaining $j$ is at least $A_{N+1}$
and contributes at most $1/A_{N+1}$. A sector solution requires $\sigma(e)\ge t$, which the displayed
bound then contradicts. Since the bound is decreasing in $j$ once $S_{k-j}$ falls below $t$, the
values $j\ge J$ are all excluded together.
\end{proof}

At the floor $k=K(d)$ one has $J=1$, and there phase~$1$ alone is a complete argument, which is the
case the table above uses; at $k=K(d)+2$ one has $J=3$, and there phases~$1$ and~$2$ together are
complete. Both correspondences are verified directly at $d=47058$, at $k=54$ and $k=56$ respectively.
Since only those two phases are implemented, the reachable rungs are those with $J\le3$, and $J$ is
what prices the ladder:
\[\begin{array}{lrrrrr}
d & \omega(d) & K(d) & k & J & |P\cup Q|\ \text{if excluded}\\[2pt]
42 & 3 & 62 & 62 & 1 & 68\\
42 & 3 & 62 & 64 & 3 & 70\\
42 & 3 & 62 & 66 & 5 & 72\\
47058 & 5 & 54 & 54 & 1 & 62\\
47058 & 5 & 54 & 56 & 3 & 64\\
47058 & 5 & 54 & 58 & 6 & 66\\
2214502422 & 6 & 70 & 70 & 1 & 79\\
2214502422 & 6 & 70 & 72 & 3 & 81
\end{array}\]
at $N=400$ \textup{(\texttt{code/sector\_phasecount.gp})}. The truncation is worth minimising rather
than maximising: $J$ is what must stay small, and at $d=2214502422$ with $k=72$ it equals $3$ for every
$N\ge300$ and $4$ below that, so $N=300$ is the smallest complete choice and the one that leaves the
fewest subsets to enumerate. The value of $J$ grows quickly with $k$, and
raising the truncation barely helps, since $j/A_{N+1}$ falls only as fast as $A_{N+1}$ grows: at
$d=42$ the rung $k=66$ needs $J=5$ at every $N$ from $400$ to $3000$. That is what stops the $d=42$
sector at $|P\cup Q|\ge70$, and it is a limitation of the implemented phases rather than of the
method.

The next step at $d=47058$ costs more, because the single--phase condition fails there. With
$N=400$ and $A_{N+1}=2791$ one has $S_{55}+1/A_{N+1}=1.00629$ and $S_{54}+2/A_{N+1}=1.00309$, both
above $d/d'=1.0000213$, so one and two primes beyond the truncation must each be enumerated; but
$S_{53}+3/A_{N+1}=0.99984<d/d'$, so three or more cannot occur and the two phases together are
complete. Both are empty: phase~$1$ enumerates $4.387\times10^{11}$ leaves and $88{,}472{,}925$ exact
checks, phase~$2$ a further $373{,}644$ leaves and $59{,}276$ exact checks, and neither yields an
admissible forced prime \textup{(\texttt{code/sector\_kexclude.rs}, $12{,}658$\,s on nine cores)}.
Hence $\omega(e)\ne56$. Since $d$ is even, $e=d+qd'$ is odd and $e'=dq$ is even, so
Proposition~\ref{prop:oddsector} forces $\omega(e)$ even; with $\omega(e)\ge56$ already in hand this
gives $\omega(e)\ge58$ and
\[ |P\cup Q|\;=\;1+\omega(d)+\omega(e)\;\ge\;1+5+58\;=\;64 \]
at $d=47058$. The truncation is what stops the ladder here: at $\omega(e)=57$ even four primes beyond
$A_N$ remain possible, so the next rung needs a larger $N$ rather than more time at this one.
\end{proposition}

\begin{remark}
The $d=42$ row recovers \cite{Bado2026b}'s $\omega(e)\ge64$ by a route with nothing in common with
its proof: that one maximises mass over a $51{,}281$-state residue automaton, this one enumerates
$175$ prime sets and tests an integrality condition. Two independent derivations of the same bound.

Running the same exclusion across sectors rather than up the ladder is what the identity really
buys. Among all sectors $d$ that are squarefree with $\omega(d)\le12$ and every prime factor at most
$71$, exactly $16{,}234$ have mass floor $|P\cup Q|=60$; in every one of them the floor value of
$\omega(e)$ is excluded, the largest enumeration being $41{,}654$ sets
\textup{(\texttt{code/sector\_at60.gp}, \texttt{code/sector\_kexclude.rs})}. The counts $250$, $1{,}533$
and $16{,}234$ are reproduced from scratch by \texttt{code/sector\_at60\_count.rs}. Determining them
requires the complete factorisation of $d'$, not merely its small factors: a prime of $d'$ that goes
unfound is never matched in the exclusion set, so the primes available to $e$ are overcounted, $K(d)$
comes out too small and the floor too low, which is the direction in which a sector at $60$ could be
misread as lying below $60$ and skipped. Since $d'$ passes $10^{14}$ already at $\omega(d)=10$ and
routinely carries a large prime factor, the enumerator factors it with Pollard rho under a
Miller--Rabin test, and refuses outright when $d'$ would exceed the width its modular arithmetic is
exact to. So
\[ |P\cup Q|=60\ \text{is impossible in every such sector,} \]
and any support of size $60$ must have $\omega(d)\ge13$ or a sector prime exceeding $71$. This does
not move Theorem~\ref{K-thm:barrier}, whose $60$ is proved by a different route and over all supports;
what it does is remove the entire low--$\omega$, small--prime part of the sector decomposition from
the level--$60$ frontier.

The sectors that could reach $60$ form a \emph{finite} set, which is worth recording because the sector decomposition appears to range over unboundedly many $d$. At level $60$ the
support is a $59$--prime base together with one tail prime, and a base element $x$ satisfies
$T_{58}+1/x\ge2$, so $x\le1/(2-T_{58})=793.6\ldots$; every prime up to $167$ lies in the base, and
$d$ is a subset of the $138$ primes at most $787$. So the question ``which sectors admit
$|P\cup Q|=60$'' is a finite one. It is not, however, a small one: the count of sectors whose mass
floor is exactly $60$ rises from $250$ over primes at most $47$ with $\omega(d)\le8$ to $1{,}533$
over primes at most $59$ with $\omega(d)\le10$ and $16{,}234$ over primes at most $71$ with
$\omega(d)\le12$ (all cleared, each family a subset of the next), and it does not saturate: raising the
cap one step further, to primes at most $89$ at the same $\omega(d)\le12$, gives $125{,}813$, a factor
of $7.75$ for one more prime in the pool. Those are counted, not cleared, and are recorded here only as
the next term of the growth. Holding $\omega(d)\le8$ and letting the prime bound $B$ grow, the
count runs $26,\,83,\,250,\,588,\,1762,\,5172$ at $B=31,41,47,59,71,89$; holding $\omega(d)\le6$ it
runs $19,\,34,\,53,\,77,\,107,\,166,\,217$ at $B=31,41,47,59,71,89,107$. Against $\pi(B)$ the fitted
exponent grows with $\omega$, from about $2.7$ at $\omega=6$ to about $6.8$ at $\omega=8$, so the count
behaves like a binomial $\binom{\pi(B)}{\omega}$. With the true bounds $B=787$ and $\omega(d)\le58$ the
family is finite but of order $\binom{138}{58}$, and enumerating it is not possible. Clearing sectors
one at a time therefore strengthens the statement at each step and cannot reach the barrier. Finiteness here converts the sector programme into a
computation of the same order as the one it was meant to replace.

Across every sector enumerated so far the floor value of $\omega(e)$ is not attained, and that is a
falsifiable statement in its own right.

\begin{proposition}[Odd sectors keep the prime $2$ only when $\omega(d)$ is odd]\label{prop:oddsector}
Let $d$ be odd and squarefree. Then $d'\equiv\omega(d)\pmod2$. Consequently, if $\omega(d)$ is even
then $2\mid d'$, so $2\mid dd'$ and $2\notin\operatorname{supp}(e)$ by
Remark~\ref{rem:sectordprime}; the mass $d/d'$ must then be carried entirely by odd primes dividing
neither $d$ nor $d'$.
\end{proposition}
\begin{proof}
For odd squarefree $d$ each cofactor $d/p$ is a product of odd primes, hence odd, so
$d'=\sum_{p\mid d}d/p$ is a sum of $\omega(d)$ odd terms and has the parity of $\omega(d)$. Verified
with no violation on the $81{,}053$ odd squarefree $d<2\times10^5$.
\end{proof}

The cost of losing $2$ is not marginal, because $\tfrac12$ is the largest single contribution
available: at $d=15015=3\cdot5\cdot7\cdot11\cdot13$, where $\omega(d)=5$ is odd and $2$ survives,
$K(d)=127$, while at $d=255255=3\cdot5\cdot7\cdot11\cdot13\cdot17$, where $\omega(d)=6$ is even and
$2$ is lost, $K(d)=1535$. Hence a sector with $d$ odd and $\omega(d)$ even has mass floor far above
$60$, and every sector of mass floor $60$ with $d$ odd has $\omega(d)$ odd. That is what the
enumeration shows: of the $16{,}234$ sectors of mass floor $60$ over primes at most $71$, the
$4{,}846$ with $d$ odd have $\omega(d)\in\{9,11\}$ and none has $\omega(d)$ even, although nothing in
the enumeration imposes it.

\begin{conjecture}[The floor gap]\label{conj:floorgap}
For every squarefree $d$, no two--cycle in the sector $d$ has $\omega(e)$ equal to the mass floor
$K(d)$, the least $k$ with $S_k(d)\ge d/d'$.
\end{conjecture}

The evidence is $5{,}680$ sectors with no attained floor: the $1{,}533$ of mass floor $60$ over primes
at most $59$, a further $3{,}600$ of the family over primes at most $71$, the three sectors of
Proposition~\ref{prop:ppnsectors}, and $480$ sectors of mass floor $61$ through $68$ over primes at
most $47$, the last of these a range disjoint from all the others and tested after the statement was
fixed. Its consequence is that the barrier would move: the minimum of $1+\omega(d)+K(d)$ over sectors
is $60$, attained at $d=2\cdot3\cdot7\cdot23\cdot41$ and checked over $109{,}293$ sectors, so with the
conjecture and the parity of Proposition~\ref{prop:ppnsectors} every sector would give
$|P\cup Q|\ge62$. The evidence must be read with its power, which is negligible. The integrality
$\Delta_d(R)\mid d'R+d^2$ has probability of order $1/\Delta_d(R)$, and $\Delta_d(R)$ is of the size
of a product of some sixty primes; over the $5{,}680$ sectors and their some $10^4$ candidate sets
each, the expected number of attained floors is of order $10^{-92}$ whether the conjecture is true or
false. An enumeration of this kind therefore cannot support it, and no feasible enlargement would.
Nor is the floor structurally distinguished: the rate at which candidates reach the exact divisibility
test matches the width of the numerical prefilter at the floor and above it alike, the ratios of
observed to expected being $0.72$ and $0.86$ at the floor for $d=47058$ and $d=2214502422$ and $1.01$
at $\omega(e)=64$ for $d=42$, so the deficit at the floor is the small number of candidate sets and
nothing else. The conjecture is recorded because it is what every enumeration performed here
verifies and because its consequence is sharp, not because the enumerations are evidence for it.

The ladder does not continue cheaply. In every sector the first rung above the mass floor is small, the three counts above, and the next requires the two-phase machinery of
Proposition~\ref{prop:sector42}, because the truncation bound no longer excludes primes beyond
$A_N$: at $d=42$ that step cost $3.7\times10^{10}$ sets, and at $d=47058$ the analogous
$\omega(e)=56$ is $4.4\times10^{11}$, not run. The cheapness of the first rung and the cost of the
second is the shape of the method, not an accident of $d=42$.
\end{remark}

\begin{remark}
The same argument decides $\omega(e)=66$ in principle, but with up to four primes beyond the
truncation the nested ranges multiply, and the enumeration is no longer a twenty--minute job; we
have not run it. What the proposition shows is where the sector programme actually stops: not at
the residue sieve, whose limit is $68$, but at the cost of an exact enumeration in which the last
one or two primes are forced by the rest. The two--large--prime case, a factoring problem in general
($(\Delta\ell_1-41R)(\Delta\ell_2-41R)=1681R^2+1764\Delta$), is finite here only because the smaller
of the two is trapped in $(1/\delta,2/\delta)$ and $\delta$ is bounded below over a finite list.
\end{remark}

Three consequences, in increasing order of what they cost.

\emph{The gate is finite, not Lehmer.} The tail primes of a base lie in the explicit finite set
$\{2(N\pm k)/m^2: m\mid 4D,\ k^2=AB+Dm^2\}$. So the second of the three gates is a search over the
divisors of $4D$, not a primality question along an exponential sequence, and the classification of
the residual difficulty as Lehmer--class rather than Schinzel--class does not survive. Checked sound
\emph{and complete} against exhaustive search on $1200$ bases with zero mismatches, and stepped out
along a real orbit, in \texttt{code/tailbound.gp}: on the base $D=1$, $N=7$ (bound $N=7$) the orbit
gives $q_1=7$ prime, then $q_2=741895$, $q_3=76921173511$, and every later term composite, with both
square conditions holding throughout.

\emph{The count reverses sign, and the search has now been run.} The expected number of admissible
$m$ is $\sum_{m\mid 4D}(AB+Dm^2)^{-1/2}\approx D^{-1/2}\sum_{m\mid4D}m^{-1}$, which on the first
immune base is $2.4\times10^{-56}$ and below $10^{-38}$ summed over all $34$ immune families, the
opposite of the ``weakly favours \textsc{yes}'' of Remark~\ref{rem:killcount} on exactly the families
that reading was about. That prediction is now checked. The full divisor space is
$4\cdot2^{58}\approx1.15\times10^{18}$ per family and is not enumerable, but because the hit
probability falls like $1/(m\sqrt D)$ the expected count is carried by small $m$: restricting to
$\omega(m)\le8$ leaves $99.9998\%$ of the expected count. Over all $34$ families that slice is
$3.0777\times10^{11}$ candidates, and
\[ \textup{it contains no }m\textup{ with }AB+Dm^2\textup{ a perfect square} \]
($4{,}933$ survive the modular cascade, exact \texttt{issquare} kills all of them; the earlier
$\omega(m)\le7$ pass, $4.7081\times10^{10}$ candidates and $771$ survivors, agrees).

The same search has been run over \emph{every} admissible base, which matters because
Proposition~\ref{K-prop:tailbound} requires no primality of $A_S$ or $B_S$, only the algebra. The $34$
immune families are the ones left after the Jacobi kill and the factoring attack, and their immunity
is conditional on a Baillie--PSW verdict; the other $3{,}444$ of the $3{,}478$ open families are open
precisely because $A_S$ or $B_S$ is a balanced semiprime nobody can split. The tail search does not
care. Over all $49{,}961$ admissible bases at $\omega(m)\le4$, which is $99.2\%$ of the expected count
per base, the run is $9.1296\times10^{10}$ candidates, $1{,}469$ survive the modular cascade, and
exact \texttt{issquare} kills all $1{,}469$. Every survivor was additionally checked to satisfy
$m\mid 4D$, the hypothesis the proposition needs. So
\[ \textup{no admissible base admits a Pythagorean tail prime with }\omega(m)\le4, \]
unconditionally and with no primality input. Where Proposition~\ref{prop:tailkill} leaves $3{,}478$
of the one--new--prime families at level $60$ open, of which $34$ only conditionally immune, this
leaves none open on the stated slice. It is a slice result and not a proof: the untested remainder is
$0.8\%$ of the expected count per base, and $\omega(m)\ge5$ is untested.
The run is \texttt{code/tailsearch.rs} (a cascade of quadratic--residue filters modulo primes coprime
to $2D$, $771$ survivors) with every survivor then tested exactly by \texttt{issquare} in
\texttt{code/tailsearch\_verify.gp} ($0$ squares). A positive control recovers the known solutions of
the toy base $D=1$, $N=7$, and an exhaustive exact run at $\omega(m)\le2$ agrees with the filter, so
the cascade is not silently over--rejecting.

It is worth recording what the slice covers in the other measure, since the two are thirteen orders
apart. The parametrisation runs over $m\mid4D$ with $m$ even, so $m=2\alpha$ with $\alpha\mid D$, and
at level $59$ that is $2^{59}=5.7646\times10^{17}$ divisors. Splitting on the parity of $\alpha$, the
divisors with $\omega(m)\le4$ number $2\sum_{j\le3}\binom{58}{j}=65{,}136$, and those with
$\omega(m)\le8$ number $2\sum_{j\le7}\binom{58}{j}=692{,}376{,}800$. As fractions of the space the
bound $q\le N$ actually proves finite, these are
\[ 1.130\times10^{-13}\qquad\text{and}\qquad 1.201\times10^{-9}, \]
against the $99.2\%$ and $99.9998\%$ of the \emph{expected count} quoted above. Both figures are
correct and they measure different things: the expected count concentrates on small $\omega$, the hit
probability falling like $1/(m\sqrt D)$, so a slice can carry almost all of the expectation while
carrying almost none of the space. Only the second pair of numbers is a statement about what has
been checked. This is \textsc{verified} on the stated slice, not on the
full divisor set: the residual $0.0018\%$ of the expected count is untested.

The filter primes must be coprime to $2D$, which is itself a structural fact rather than a
convenience. Since $D$ contains every prime up to $167$, for any $\ell\mid D$
\[ AB+Dm^2-N^2=D\bigl(m^2-4D\bigr)\equiv0 \pmod{\ell}, \]
so the value is the square $N^2$ modulo $\ell$ for \emph{every} $m$, and as $\gcd(D,N)=1$ Hensel
lifts that root to every power of $\ell$. The tail search therefore has no local obstruction at any
prime dividing $D$, the analogue of Proposition~\ref{prop:localcomplete} for this formulation; the
classical $64/63/65/11$ pre--filter for squares rejects nothing here, measured at a $75\%$ pass rate
modulo $64$ and $100\%$ modulo $63$. Formalised as \texttt{tail\_value\_mod\_dvd} and
\texttt{hensel\_lift\_identity}; the induction from the step to all powers is routine and is not
formalised.

\emph{The paper's own height estimate now says the same thing.} Remark~\ref{rem:lehmer} puts the
least point of a nonempty fiber at height $\exp\bigl(10^{112}\bigr)$ on regulator grounds, so its
$q\approx x^2/A$ carries some $10^{112}$ digits, against a bound of $115$. Two independent estimates,
one from divisor counting and one from the class--number formula, agree that no immune family admits
a Pythagorean tail prime. Neither is a proof: the bound $q\le N$ is proved, the emptiness is
\textsc{heuristic}.

Formalised in \texttt{lean/Erdos307/TailBound.lean} (\texttt{tail\_factor}, \texttt{tail\_quadratic},
\texttt{tail\_discriminant}, \texttt{tail\_bound}, \texttt{tail\_finite}) on the three standard
axioms, with a non--vacuity witness at the base above.
\end{remark}

\begin{remark}[The immune signal and its Pell orbit]\label{rem:lehmer}
The orbit below is real, but primality along it is not the residual question, since
Proposition~\ref{K-prop:tailbound} bounds the tail prime by $N$.

For a reciprocity--immune base the Pell system of Remark~\ref{rem:tier60},
$B_Sx^2-A_Sy^2=-4D_S^2$, is locally soluble everywhere ($\kappa_S>0$), and a Pythagorean pair in the
tail family $S\cup\{q\}$ is \emph{exactly} a solution with $q=(x^2-D_S)/A_S=(y^2-D_S)/B_S$ a
positive prime. Its solutions grow like $\varepsilon^n$ in the fundamental unit of the order of
discriminant $4A_SB_S$ ($224$ digits, $A_SB_S$ a nonsquare, on an explicit immune base). Writing
$x_{n+1}=tx_n-x_{n-1}$ with $x_n=(\alpha^n+\beta^n)/2$ and $\alpha\beta=1$, the squares $x_n^2$ have
characteristic roots $\alpha^2,\beta^2,1$ and the constant shift is absorbed by the root $1$, so
\[ q_n=\frac{x_n^2-D_S}{A_S}\qquad\text{satisfies}\qquad q_{n+3}=(t^2-1)q_{n+2}-(t^2-1)q_{n+1}+q_n, \]
a non--degenerate ternary linear recurrence (verified on four Pell parameters). A
reciprocity--killed base is the complementary degeneration: there the same system is locally
\emph{insoluble} ($\kappa_S=0$), so no candidate exists at all, and the dichotomy of
Proposition~\ref{prop:tailkill} is Pell--solubility versus Pell--insolubility.

The return is milder than ``$k=0$'' suggests: on the Pythagorean locus the deviation $k=(a'-b)/a$ is
an \emph{integer} (the identity $v(D(u)-v)=u(u-D(v))$ with $\gcd(u,v)=1$), forced into $\{0,-1\}$ by
the near--balanced immune ratio, so $k=0$ is a sign choice rather than a measure--zero coincidence,
and the naive Wieferich shortcut that would pin $q$ is vacuous, since $x^2\equiv D_S\pmod q$ holds
identically from $x^2=A_Sq+D_S$.

Neither square condition carries the difficulty alone. Imposing only the minus--square,
$d^2=B_Sq+D_S$, the admissible $d$ satisfy $d^2\equiv D_S\pmod{B_S}$, and along such a class
$d=d_0+tB_S$ one has
\[ q(t)\;=\;\frac{d_0^2-D_S}{B_S}\;+\;2d_0\,t\;+\;B_S\,t^2, \]
a quadratic in $t$ (toy bases: $3+54t+173t^2$, $7+238t+1693t^2$, $2+446t+17357t^2$): a Bunyakovsky
question, of the same class as the $p=5$ family of Proposition~\ref{C-prop:plus}\textup{(iii)}, and
conjecturally routine. It is the \emph{conjunction} of the two squares that replaces the quadratic
parameter by a fundamental unit and turns the orbit exponential.

\emph{Where the reading broke.} That exponential growth was taken to place the question in the
Lehmer class, on two grounds: no algebraic factorisation is available, $D_S$ being squarefree so
that $x^2-D_S$ is irreducible, and no covering congruence can kill the family, by
Proposition~\ref{prop:localcomplete}. Both observations are correct, and neither is the relevant
one. The factorisation that does exist is not of $x^2-D_S$ but of $x^2-y^2=4D_Sq$, available because
$A_S-B_S=4D_S$ \emph{exactly}: the identity was used here only to build the Pell system, never
multiplicatively. Proposition~\ref{K-prop:tailbound} draws it, and the ternary recurrence above then
supplies no prime term past $q\le N$.

Three nested gates separate an immune family from a solution: fiber nonemptiness, a tail prime, and
the $k=0$ return. The first is a class--group condition, since Hasse supplies only rational points,
and it is quantitative: immunity makes $A_SB_S$ a product of two distinct primes, hence squarefree,
so the conic lives in the near--maximal order of $\mathbb{Q}(\sqrt{A_SB_S})$ of discriminant
${\sim}A_SB_S\approx10^{224}$, and the class--number formula puts the expected regulator at
$R=\sqrt\Delta\,L(1,\chi)/(2h)\approx10^{112}$; barring an anomalously large class number, the
\emph{least} point of a nonempty fiber sits at height $\exp\bigl(10^{112}\bigr)$. That estimate is
used in Remark~\ref{rem:notlehmer} as independent corroboration that the second gate is empty. A
direct sweep to $q\le10^{12}$ was structurally incapable of reaching even the first gate.
\end{remark}

\begin{remark}[The affine sieve does not reach a rank--one orbit]\label{rem:affinesieve}
The modern instrument for primes in orbits is the affine sieve of Bourgain, Gamburd and Sarnak
\cite{BGS2010}, which produces $r$--almost--primes in orbits of thin groups. It is the one method
designed for exactly the shape of question Remark~\ref{rem:lehmer} isolates, so its failure here is
worth recording precisely rather than left implicit.

The affine sieve needs a spectral gap: the congruence quotients of the acting group must form an
expander family, and that is what supplies the level of distribution the sieve consumes. Here the
$q_n$ run along a Pell orbit, so the acting group is the cyclic group generated by the fundamental
unit of the order of discriminant $4A_SB_S$. Cyclic groups are abelian, hence amenable, and have no
spectral gap; expansion is precisely the property an abelian group lacks. The sieve therefore
returns nothing, and it returns nothing on $2^n-1$ for the same reason. The obstruction is rank, not
thinness as such: non--elementary thin orbits, Apollonian packings among them, do yield
almost--primes.

The natural repair fails for a reason worth stating too. One may embed the rank--one orbit in a
larger thin group whose orbit does expand, and sieve there. What the affine sieve then delivers is
infinitely many almost--primes in the \emph{larger} orbit, inside which our sequence has relative
density zero; the sieve does not localise to a prescribed sub--orbit. Expansion is bought at the
cost of losing the sequence.

This meets the unit--rank trichotomy of Remark~\ref{B-rem:dichotomy} from a second direction. There,
rank zero gave a classical gate and denied a continuous family, positive rank the reverse. Here,
rank one is exactly the cell in which the affine sieve has no traction. Stated as a single missing
property: what is wanted is a substitute for expansion at rank one, and that is the Mersenne problem
itself, not a lemma about \#307.
\end{remark}

\begin{remark}[The gap to the record, measured]\label{rem:sparserecord}
It is worth stating how far the required statement is from the best that is known, because the
distance is not one of degree.

Write a sequence's \emph{density exponent} $\theta$ for $\#\{n\le X\}\asymp X^{\theta+o(1)}$. The
sparsest sets currently proved to contain infinitely many primes are the values of binary forms:
$a^2+b^4$ at $\theta=3/4$, by Friedlander and Iwaniec \cite{FI1998}, and $a^3+2b^3$ at $\theta=2/3$,
by Heath-Brown \cite{HB2001}, the latter being the sparsest polynomial form known to capture its
primes. Our $q_n$ grow like $\varepsilon^{2n}$, so $\#\{n:q_n\le X\}\asymp\log X$ and
\[ \theta=0. \]

The distance from $2/3$ to $0$ is categorical rather than quantitative. Both records are bilinear
arguments: they run a sieve of dimension one and close it with a Type~II estimate obtained by
averaging over a set of size a positive power of $X$. Every gain in such an argument is a gain of a
power of $X$, and against $\log X$ terms a power of $X$ gains nothing; below $\theta=1/2$ the
standard sieve axioms already fail, and at $\theta=0$ there is no set left to average over. This is
the same wall as Remark~\ref{rem:affinesieve} approached from the analytic rather than the spectral
side, and it is why every sequence of exponent zero, Mersenne, Fibonacci, Lucas and ours alike, is
open together.

The structure our sequence does carry is the standard one and confers nothing. From
$A_Sq_n=x_n^2-D_S$ one gets $x_n^2\equiv D_S$ for every prime $p\mid q_n$ not dividing $D_S$, so all
prime factors of $q_n$ lie in the density--$\tfrac12$ set where $D_S$ is a quadratic residue. That
constrains the factorisations without excluding any, which is exactly the position the primitive
divisor theorems already left us in: divisors, never primality.
\end{remark}

\begin{remark}[The residue, atomised, and where its difficulty comes from]\label{rem:atomised}
Decomposing the question into independently known statements locates the obstruction on a single
atom and identifies its source.

Two atoms carry theorems. Every $q_n$ beyond the thirtieth has a primitive prime divisor
\cite{BHV2001}; and the greatest prime factor satisfies
$P^+(q_n)>n\exp\bigl(\log n/(104\log\log n)\bigr)$, by Stewart \cite{Stewart2013}, which settled
questions of Schinzel and Erd\H{o}s and was the first general advance on Bang and Carmichael. Every
other atom, that $q_n$ be squarefree infinitely often, that $\omega(q_n)$ be bounded infinitely
often, that $q_n$ be a $P_k$, and primality itself, is open.

The chain that would close the crux is $P^+$ large, hence $\omega$ small, hence prime. It breaks at
the first link, and by a measurable amount. Stewart's bound certifies
$\log P^+\ge(1+o(1))\log n$ against $\log q_n\approx2n\log\varepsilon$, so the fraction of the
term's size that is certified is of order $\log n/n$: $10.6\%$ at $n=10$, $0.32\%$ at $n=10^3$,
$4\times10^{-4}\%$ at $n=10^6$. Primality is the fraction $1$. No sharpening of the constant $104$
changes this, since $n^{1+o(1)}$ is the shape linear forms in logarithms produce. The Subspace
Theorem is the other available instrument and does not help: it is stronger under weaker hypotheses
but ineffective \cite{BCZ2003}, and an ineffective statement cannot exhibit the single prime term a
construction needs.

The source is visible in the comparison with the polynomial case. For $n^2+1$ the greatest prime
factor exceeds $n^{1.3}$ infinitely often \cite{Merikoski2019}, that is $q^{0.65}$, a \emph{positive
proportion} of the term's size, and on such sequences almost-primality is available too. For an
exponentially parametrised sequence only $(\log q)^{1+o(1)}$ is available, a vanishing proportion.
The gap between the two regimes is exactly the gap between polynomial and exponential
parametrisation, and Proposition~\ref{K-prop:nopoly} proves \#307 admits no polynomial one. So the
obstruction recorded in Remark~\ref{rem:affinesieve} as a question of rank, in
Remark~\ref{rem:sparserecord} as one of density exponent, and here as one of $P^+$, is a single
obstruction with three faces, and its source is Proposition~\ref{K-prop:nopoly}: forced exponentiality
is what places the residue in the class where nothing of the required strength is known.

Two things sharpen that attribution, and both make the obstruction less incidental rather than more.
First, the boundary is exactly where Proposition~\ref{K-prop:nopoly} sits. The fraction of a term's
size that Stewart's bound certifies is $\log P^+/\log q_n\approx\log n/\log q_n$, so it stays bounded
below precisely when $\log q_n\ll\log n$, that is when $q_n$ grows polynomially: degree $d$ gives the
constant $1/d$. Every super-polynomial rate sends it to $0$, and not only the exponential one:
$\exp(\sqrt n)$ and $n^{\log n}$ do as well. So polynomial growth is not one convenient case among
several, it is the whole of the favourable region, and it is exactly what
Proposition~\ref{K-prop:nopoly} excludes.

Second, the hypotheses of Proposition~\ref{K-prop:nopoly} are worth testing, since a gap there would
dissolve everything above. It assumes a \emph{fixed} number of polynomials, and its proof uses that
finiteness when it splits $\sum_i1/p_i$ into constants plus a remainder tending to $0$. A family
whose support grows with the parameter is therefore not covered. That gap is real and useless: a
growing support means a growing number of simultaneous primality conditions, which is a harder
demand than one, not an easier one, and it does not return the family to the polynomial regime. The
proposition forbids the case that would help and leaves uncovered only cases that would not.
\end{remark}

\begin{remark}[The squarefree atom, and what it reduces to]\label{rem:squarefree}
Among the four open atoms one is structurally easier, and it is worth saying why and how far it
gets. Squarefreeness of $q_n$ asks only for an \emph{upper} bound on a count of bad events, never
for the production of a prime, so it is the one atom not subject to the parity obstruction that
blocks bounded $\omega$, $P_k$, and primality.

The argument it wants is a union bound. For a prime $p\nmid2A_SD_S$, $p^2\mid q_n$ forces
$x_n^2\equiv D_S\pmod{p^2}$; the Pell orbit is purely periodic modulo $p^2$ with period
$\pi(p^2)$, there are at most two square roots of $D_S$, and each is met a bounded number of times
per period, so the density of bad $n$ attached to $p$ is at most $C/\pi(p^2)$. Where
$\pi(p^2)=p\,\pi(p)$, which is the generic behaviour, and $\pi(p)\asymp p$, this is $\asymp C/p^2$,
and the union bound closes as soon as $C\sum_pp^{-2}<1$. The constant is comfortable:
$\sum_pp^{-2}=0.4522474\ldots$, so any $C<2.21$ suffices, and a positive proportion of the $q_n$
would be squarefree, hence infinitely many.

What is missing is the hypothesis $\pi(p^2)=p\,\pi(p)$. Its failure is precisely a
Wall--Sun--Sun prime for the sequence, and there the density is $\asymp1/p$ rather than
$\asymp1/p^2$, whose sum diverges and destroys the bound. No such prime is known for any classical
sequence and their finiteness is unproved, so the atom does not close: it \emph{reduces} to a
Wall--Sun--Sun statement. That is a sharper position than ``open'', since it names the obstruction
and shows it is a single hypothesis rather than a programme, but it is a reduction and not a proof.

The unconditional pieces are formalised. \texttt{Erdos307.primeSquareSum\_certificate} gives
$\sum_pp^{-2}<\tfrac12$ in the finite form used above, the first six primes together with the bound
$\tfrac12\sum_{m\ge16}m^{-2}<\tfrac1{30}$ for the odd primes beyond; \texttt{good\_density\_pos} and
\texttt{squarefree\_positive\_density\_of\_bound} are the union bound and the conditional
implication. All carry $0$ \texttt{sorry} and the three standard axioms
\textup{(\texttt{lean/Erdos307/Squarefree.lean})}.

Two further points fix the standing of this atom, and both are negative. First, on the immune bases
the argument is blocked at its \emph{lowest} rung as well as its highest. The density attached to a
single small $p$ is $\#\{n\bmod\pi(p^2):p^2\mid q_n\}/\pi(p^2)$, read off the Pell orbit modulo
$p^2$, and that needs the fundamental solution, the object Remark~\ref{B-rem:untestable}\,(i) places
at $10^{24.7}$ operations. For a Lucas sequence with a known fundamental unit the small-prime rung is
a finite computation; here it is not. The ladder is therefore blocked at both ends, by
non-computability below and by Wall--Sun--Sun above. Second, and more to the point, squarefreeness is
not what the construction needs: Remark~\ref{rem:lehmer} requires $q$ to be \emph{prime}, and a
squarefree $q$ with several prime factors is useless there, so even the dream lemma of this atom
would leave \#307 exactly where it stands. Of the four atoms this is the easiest, and it is also the
only one whose resolution would not advance the problem.
\end{remark}

\section{Local anatomy and a heuristic model}\label{sec:anatomy}

\begin{lemma}[Local anatomy]\label{lem:anatomy}
Let $a$ be squarefree and $q$ a prime with $q\nmid a$. Then $q\mid a'$ if and only if
$\sum_{p\mid a}p^{-1}\equiv 0\pmod q$, where $p^{-1}$ is the inverse of $p$ modulo $q$. If $q\mid a$,
then $q\nmid a'$.

\textup{Formalised.} \texttt{Erdos307.anatomy\_iff}, with the factorisation
$a'\equiv a\sum_{p\mid a}p^{-1}$ proved as \texttt{csum\_eq\_mul\_recipSum} rather than assumed; the
second half is \texttt{anatomy\_on\_support} and is Rigidity
\textup{(\texttt{lean/Erdos307/Frame.lean})}.
\end{lemma}
\begin{proof}
$a'=a\sum_{p\mid a}1/p$. If $q\nmid a$ then $q\mid a'$ iff $q\mid a\sum_{p\mid a}p^{-1}$ iff
$\sum_{p\mid a}p^{-1}\equiv0\pmod q$. If $q\mid a$, then by Lemma~\ref{K-thm:rigidity}
applied to the prime set of $a$, $q\nmid a'$.
\end{proof}

Thus the cycle equation $a'=b$, $b'=a$ factors, prime by prime, into
$\omega(a)+\omega(b)$ local congruences of exactly the shape that governs amicable--pair heuristics:
for each $q\mid b$ one needs $\sum_{p\mid a}p^{-1}\equiv0\pmod q$, and symmetrically. These local
conditions, however, are individually \emph{empty of obstruction}:

\begin{proposition}[No single--modulus congruence obstruction]\label{prop:nolocal}
For every modulus $m$ with $2\le m\le210$, the reduced cross--match system $N_P\equiv D_Q$,
$N_Q\equiv D_P\pmod m$ is realisable by disjoint squarefree prime sets $P,Q$, in both structural
branches: $\gcd(m,D_PD_Q)=1$, and the branch in which a prime divisor of $m$ lies in $P\cup Q$. Hence
no individual modulus in this range can rule out a solution of \#307.
\end{proposition}
\begin{proof}[Construction]
Write $\Sigma_S\equiv\sum_{p\in S}p^{-1}\pmod m$, so $N_S\equiv D_S\Sigma_S$. The reduced system is
$\Sigma_P\equiv D_Q D_P^{-1}$, $\Sigma_Q\equiv D_P D_Q^{-1}\pmod m$; its only coupling is
$\Sigma_P\Sigma_Q\equiv1$, automatic, with $D_P,D_Q$ otherwise free. By Dirichlet each unit class mod
$m$ contains primes, fixing $D_P,D_Q$; appending a prime $r\equiv1\pmod m$ leaves $D_S$ fixed while
sending $N_S\mapsto rN_S+D_S\equiv N_S+D_S$, i.e.\ incrementing $\Sigma_S$ by $1$, so (when
$\gcd(D_S,m)=1$) every target class is reached with disjoint sets. In the branch $\ell\mid m$,
$\ell\in P$, reduction mod $\ell$ collapses the system to $N_Q\equiv D_P\equiv0\pmod\ell$, met by
inverse--paired residues in $Q$; composite $m$ needs no separate gluing, the increment acting directly
mod $m$. We verified both branches for all $2\le m\le210$ (for prime and prime--power moduli the
inside branch requires the directed construction above; seeding a prime divisor of $m$ into $P$,
a blind search missing these as a search--power artefact, not an arithmetic obstruction).
\end{proof}
\begin{remark}[Precise scope]
This excludes only a \emph{single fixed--modulus} congruence disproof. The witnesses are
$m$--dependent (the modulus--$6$ witness already fails modulo $210$), and a single $(P,Q)$ solving the
system for \emph{all} $m$ simultaneously is, for bounded integers, equivalent to the integer identities
$N_P=D_Q$, $N_Q=D_P$ (that is, to \#307 itself) so per--modulus solvability is \emph{no} evidence
for existence. Nor does it exclude other finite or local negative arguments (size/Mertens bounds,
density, magnitude--congruence hybrids, or a Hasse--type failure across infinitely many moduli);
indeed our own $|U|\ge59$ and $2.09\times10^{56}$ barriers are finite--data results fully consistent
with it. Together with Theorem~\ref{B-thm:ff} (every non--archimedean shadow trivially negative for
degree reasons) it brackets the problem: no single--modulus obstruction over $\mathbb{Z}$, fatal
non--archimedean analogues over $\mathbb{F}_q[t]$, the difficulty living at the archimedean place.
\end{remark}

\paragraph{A first--order count.}
Model a candidate squarefree $a$ as ``$P$--abundant'' ($\sum_{p\mid a}1/p>1$) and ask for
$b=a'$ to itself be $Q$--abundant with $\sum_{q\mid b}1/q=1/s$. Two empirical inputs (Section
\ref{K-sec:computations}) calibrate this:
\begin{itemize}[leftmargin=1.4em]
\item the density of $P$--abundant $n$ is $\delta=0.0420$ (and $\approx0.0176$ among squarefree
$n$); a theorem, not merely an observation (Proposition~\ref{prop:density});
\item for a qualifying squarefree $a$, $b=a'$ is squarefree about $95\%$ of the time, but its mass
$\sum_{q\mid b}1/q$ has mean $\approx0.047$ against the required $\approx0.934$; the empirical tail
probability $\Pr[\sum_{q\mid b}1/q\ge 1/s]$ is $0$ in $40{,}920$ samples at these scales
(Section~\ref{K-sec:computations} of the core note).
\end{itemize}
The gap is structural: at small scale a $P$--abundant $a$ hogs the cheap mass--carrying small primes,
and $b=a'$, being coprime to $a$, is both too small and too coprime to assemble reciprocal mass
$1/s$. The extremal case is stark: for the primorial $a=\prod_{i\le k}p_i$, though $\sum_{p\mid a}1/p$
rises through $2$ at $k=59$, the cofactor sum $a'$ is essentially prime, with $\sum_{q\mid a'}1/q<0.2$
for every $k<100$ (and only $0.0014$ at $k=59$): the family that \emph{maximises} the forward mass is
the most starved on the return. Equivalently, for squarefree $n$ with $n'$ squarefree,
$n''=n\,\sigma(n)\,\sigma(n')$, so the two--step return ratio $n''/n=\sigma(n)\sigma(n')$ is the \#307
product itself; over \emph{all} such $n<10^6$ it stays strictly below $1$, never exceeding $0.54$
(maximal at $n=969969=3\cdot7\cdot11\cdot13\cdot17\cdot19$), so $D$ is two--step contracting throughout
the tested range and must reverse to unit ratio for a solution to exist. The first scale at which $b$
could plausibly carry the required mass is of the order of the
barrier ($\sim10^{56}$); this dynamical mass--gap is \emph{consistent with} the proved barrier rather
than an independent derivation of it (the empirical tail probability $0$ at small scale is equally
consistent with \textsc{no} and with \textsc{yes}--beyond--the--barrier). A naive ``expected number
of solutions'' built from the local congruences (one factor $\sim1/q$ per condition) diverges like a
constant times $\log X$; this is a divergent naive count of exactly the type known to be unreliable,
so we treat it only as very weak, non--rigorous suggestive evidence that solutions could exist at or
beyond the barrier.

\begin{remark}[Calibration]
We present the count of the previous paragraph as a heuristic, not a theorem. Heuristics of
amicable--pair type are known to be delicate; we report it as weak evidence, consistent with the
divergence of the refined expected count, that the answer to \#307 is plausibly \textsc{yes}
(and hence that the period--one form of the Ufnarovski--\AA hlander conjecture is plausibly false).
\end{remark}

\begin{remark}[A cautionary refinement]\label{rem:refined}
One can try to sharpen this into a quantitative rate using a Kubilius/Erd\H{o}s--Wintner model,
stratified over which small primes lie in the support of $a$ versus remain available to $b$. Two
lessons emerge, both methodological. First, a plain Monte--Carlo estimate \emph{undercounts} the rate
by some $10^{61}$: the expectation is dominated by rare, balanced splits of the small primes that
sampling almost never sees, so a naive estimate of this kind is misleadingly small and only an
explicit stratified sum corrects it. Second, and this is why we record no numerical prediction,
the stratified estimate is itself unstable: refining the stratification from the first $8$ to the
first $14$ primes drops the rate by some $17$ orders of magnitude and pushes the implied first
occurrence from $\sim10^{10^{38}}$ out past $\sim10^{10^{55}}$, with no sign of convergence. The only
robust conclusion is negative: every version places a hypothetical first solution \emph{doubly
exponentially} beyond any conceivable search, consistent with the constructive deficit of
Section~\ref{sec:breeder}; the model does not reliably decide existence in either direction.
\end{remark}

\begin{proposition}[The abundance density is a theorem]\label{prop:density}
The additive function $f(n)=\sum_{p\mid n}1/p$ satisfies the Erd\H{o}s--Wintner three--series
criterion \cite{Tenenbaum} (here $\sum 1/p^2$ and $\sum 1/p^3$ converge), so $f$ has a limiting
distribution and the
density $\delta=\lim_{x\to\infty}\tfrac1x\#\{n\le x:f(n)>1\}$ exists, equal to $\Pr\bigl[\sum_p
X_p/p>1\bigr]$ for independent $X_p\sim\mathrm{Bernoulli}(1/p)$. Moreover the constant carries a
\emph{certified} enclosure:
\[ 0.04202\;\le\;\delta\;\le\;0.04223. \]
Writing $S_1=\sum_{p\le P_1}X_p/p$, $P_1=3\times10^4$, and $T\ge0$ for the tail, monotone upper and
lower cdf envelopes of $S_1$ are propagated through the Bernoulli convolution on a grid of step
$10^{-7}$ with the shift $1/p$ rounded down for the upper and up for the lower envelope (which
preserves them, cdfs being nondecreasing) so $\delta\ge1-U(1)$ costs nothing (the tail is
nonnegative), while $\delta\le1-L(1-a)+\Pr[T\ge a]$ with the explicit Chernoff bound
$\Pr[T\ge a]\le\exp(1.36-P_1a)$, $a=10^{-3}$ (from
$\log\mathbb{E}\,e^{\lambda T}\le\lambda s_2+0.72\lambda^2s_3$ at $\lambda=P_1$, using
$s_2<1/P_1$, $s_3<1/(2P_1^2)$ and $e^x-1\le x+0.72x^2$ on $[0,1]$); float slack $10^{-9}$ exceeds
the accumulated rounding by nine orders, and the machinery reproduces an exact four--prime
reference to $2\times10^{-16}$ (\texttt{code/delta\_cert.py}). The point value $0.0420$, at the
bottom of the enclosure, agrees to four places with the sieve to $2\times10^7$; the upper edge is
conservative, carried almost entirely by the $a$--window. As for the density of abundant integers,
no closed form is expected.
\end{proposition}

\begin{remark}[Two dynamical probes: no second obstruction]\label{rem:dynamical}
Two measurements sharpen the picture beyond the first--order mass count, and both cut towards rather
than against a solution. \emph{(i) Image density.} A two--cycle needs \emph{both} members in the image
$D(\mathbb{N})$ (each is the other's derivative), so one might fear a second obstruction; that the
high--mass numbers a cycle requires are seldom derivatives. The opposite holds: of the integers below
$10^7$, $68.3\%$ lie in $D(\mathbb{N})$, and the fraction \emph{rises} with mass, from $61.7\%$ at
$\sum_{p\mid n}1/p<0.2$ to $78.2\%$ at $\sum_{p\mid n}1/p>1$ and $80\%$ beyond. The high--mass numbers
are thus over--represented among derivatives; the second cycle condition is comfortably met.
\emph{(ii) The barrier is where the balanced mass--product reaches $1$.} Set
$r(n)=\bigl(\sum_{p\mid n}1/p\bigr)\bigl(\sum_{q\mid n'}1/q\bigr)=D^2(n)/n$, so $r=1$ is exactly the
cycle condition. Exhaustively over $n\le X$ the maximum of $r$ climbs $0.50\to0.62$ as $X$ runs
$10^3\to10^9$ (with no exact cycle, as the barrier demands), with no ceiling; the climb follows the
balanced split $r\approx(S(k)/2)^2$, $S(k)=\sum_{i\le k}1/p_i$, which reaches $1$ precisely when
$S(k)=2$, i.e.\ $k\approx59$. The dynamical trajectory thus reproduces the $|U|\ge59$ barrier from the
opposite direction. Neither probe decides existence; both \emph{remove} hypothetical obstructions,
leaving the exact archimedean closure as the sole gate.
\end{remark}

\begin{remark}[Why the analytic method is starved: $N(r)\le1$]\label{rem:nostructure}
The rank--one rigidity (Section~\ref{K-sec:status} of the core note) has a sharp consequence for the circle--method
machinery that resolves the \emph{unrestricted} Egyptian--fraction problem. For a rational $r$ let
$N(r)=\#\{\,P:\sum_{p\in P}1/p=r\,\}$ count the finite prime sets with reciprocal sum $r$. Then
\[N(r)\le1:\]
by the Rigidity Lemma (Theorem~\ref{K-thm:rigidity}) the sum $\sum_{p\in P}1/p=N_P/D_P$ is already in
lowest terms, so the reduced denominator of $r$ equals $D_P=\prod P$ and $P$ is exactly its set of
prime divisors; uniquely determined ($N(r)=1$ precisely when that denominator is squarefree and the
numerator of $r$ is its cofactor sum). The circle method needs the opposite: a target hit by a positive
proportion of subsets, so a main term dominates. This is exactly what Bloom's shift to $\{p-1\}$
manufactures, $p-1$ being composite and divisor--rich, a rational is the reciprocal sum of
\emph{many} subsets of $\{1/(p-1)\}$ (already among the first fourteen shifted primes the multiplicity
reaches $3$ and grows), whence \cite{BloomShifted} represents every positive rational there. Over the
primes themselves $N(r)\le1$ leaves nothing to average: the analytic engine that reaches $\{p-h\}$ is
provably starved on $\{p\}$, and this is the method--level reason \#307 (the rank--one, exact--prime
case) lies beyond it.
\end{remark}

\section{A breeder form and the feasibility of certificate search}\label{sec:breeder}

Allowing \emph{two} free primes on one side turns the divisor trick into a cleaner bilinear form,
Euler's amicable--pair method, in the shape used by te~Riele's breeding algorithms
\cite{teRiele1984}, and lets us state precisely the strongest honest conclusion: a
\emph{conditional} reduction, which nonetheless does not become a feasible finite search.

\begin{proposition}[Bilinear breeder form]\label{prop:bde}
Fix coprime squarefree $M_0,N$ and seek a two--cycle $a=M_0rp$, $b=Nq$ with $r,p,q$ distinct primes
coprime to $M_0N$. Put $G:=NM_0-N'M_0'\ (=\Delta_0)$, $c:=M_0N'$, and $K:=GN^2+c^2$. Then the cycle
equations $a'=b$, $b'=a$ are equivalent to
\[
(Gr-c)(Gp-c)=K, \qquad q=\frac{M_0rp-N}{N'}.
\]
\end{proposition}
\begin{proof}
Expanding, $(Gr-c)(Gp-c)-K=-N'G\,(a'-b)$ (verified symbolically), so the bilinear equation is
equivalent to $a'=b$; the second cycle equation $b'=a$ then determines $q$ as stated.
\end{proof}

\begin{proposition}[The breeder is integral]\label{prop:bdeint}
Keep the notation of Proposition~\textup{\ref{prop:bde}} and let $G\ne0$. Then
\[ (Gr-c)(Gp-c)=K \iff G\,rp-c(r+p)=N^{2}, \]
and every solution of that equation has
\[ N\ \big|\ a', \qquad q=\frac{M_0rp-N}{N'}=\frac{a'}{N}\ \in\ \mathbb{Z}. \]
\end{proposition}
\begin{proof}
Expanding, $(Gr-c)(Gp-c)-c^{2}=G\bigl(G\,rp-c(r+p)\bigr)$ while $K-c^{2}=GN^{2}$, so the two
equations differ by the factor $G\ne0$. Substituting $G=NM_0-N'M_0'$, $c=M_0N'$, $a=M_0rp$ and
$a'=M_0'rp+M_0(r+p)$ into $G\,rp-c(r+p)=N^{2}$ gives
$NM_0rp-N'M_0'rp-M_0N'(r+p)=N^{2}$, that is
\[ N\,(a-N)\;=\;N'\bigl(M_0'rp+M_0(r+p)\bigr)\;=\;N'a'. \]
Hence $N\mid N'a'$, and $N$ squarefree gives $\gcd(N,N')=1$, so $N\mid a'$ and $q=a'/N$.
\end{proof}

This is worth isolating because the displayed $q=(M_0rp-N)/N'$ is a quotient by a quantity of the
size of the barrier, and one might expect a candidate to be usable only when $N'$ happens to divide
$M_0rp-N$, a further demand of probability about $1/N'$, which would leave the construction
hopeless for reasons having nothing to do with primality. There is no such demand. The expectation
$\mathbb{E}\approx\bigl(\tau(K)/G\bigr)(\ln r\ln p\ln q)^{-1}$ is therefore the correct count: the
breeder loses candidates only to primality, never to integrality. A scan of $51{,}132$ frames
(\texttt{code/breeder\_integral.gp}) produced $9$ admissible candidates, every one with integral $q$.
Formalised in \texttt{lean/Erdos307/Breeder.lean}.

Each factorisation $K=d\cdot(K/d)$ with $d\equiv-c\ (\mathrm{mod}\ G)$ proposes a candidate
$(r,p)=\bigl((d+c)/G,(K/d+c)/G\bigr)$; a solution is any such candidate whose $r,p,q$ are all prime.
A dual, partition--side form is occasionally useful:

\begin{proposition}[Subset form]\label{prop:subset}
For a support $U$ with $D=\prod_{p\in U}p$ and $\Sigma=\sum_{p\in U}D/p$, a partition $U=P\sqcup Q$
solves \#307 iff $S=P$ satisfies $A(\Sigma-A)=D^2$, where $A:=\sum_{p\in S}D/p$. Moreover
$A(\Sigma-A)-D^2=-(1-s_Ps_Q)\,D^2$.
\end{proposition}

(Both identities verified symbolically and on $2\times10^4$ random supports.) The last identity says
the ``analytic'' residual $1-s_Ps_Q$ and the ``arithmetic'' residual $A(\Sigma-A)-D^2$ are one
equation up to the factor $D^2$: the two obvious steering handles carry no independent information.

\paragraph{A conditional reduction.} Several ingredients of a construction are unconditional. Greedy
Egyptian--fraction choice gives disjoint $P,Q$ with $0<1-s_Ps_Q=O(1/T)$ (solvability to any
precision); a prime $r$ nearest $M_0N'/\Delta_0$ gives $G\le\Delta_0^{0.475}(M_0N')^{0.525}$
unconditionally (Baker--Harman--Pintz \cite{BHP2001}), and each small prime $\ell$ can be forced to
divide $K$ by one congruence on a free prime slot of $N$ (Dirichlet then supplies the prime).
Granting these, \#307 reduces to a single prime--density statement,
\[
(\mathrm{H})\qquad\text{among the admissible candidates so produced, some triple $(r,p,q)$ is prime,}
\]
and $(\mathrm{H})\Rightarrow$ \#307 is \textsc{yes}. This $(\mathrm{H})$ is a
Hardy--Littlewood/Bateman--Horn--type assertion for an explicit, computable, \emph{non--polynomial}
family; we know of no named conjecture implying it. That is not an accident of our construction:

\begin{proposition}[No fixed--shape family, of any growth rate]\label{prop:nofamily}
Let $P(t)=\{p_1(t),\dots,p_k(t)\}$ and $Q(t)=\{q_1(t),\dots,q_l(t)\}$ be families of integers of
\emph{any} shape (no polynomial, algebraic or analytic hypothesis is imposed) indexed by
$t\in\mathbb{N}$, with $k,l$ fixed. Suppose each member takes positive prime values, each nonconstant
member tends to $+\infty$, and
$\bigl(\sum_i1/p_i(t)\bigr)\bigl(\sum_j1/q_j(t)\bigr)=1$ for \emph{infinitely many} $t$. Then every
member is constant.
\end{proposition}
\begin{proof}
Let $\varphi$ enumerate, in increasing order, the $t$ at which the relation holds; $\varphi$ is
strictly monotone, so $\varphi(n)\to\infty$. Split each side into its constant members, contributing
$A_P$ and $A_Q$, and its nonconstant ones, contributing $\varepsilon_P(t),\varepsilon_Q(t)>0$ which
tend to $0$ because each nonconstant member tends to $+\infty$ and there are finitely many. Along
$\varphi$ the relation holds at every index and the remainders still tend to $0$, so
Proposition~\ref{K-prop:nopoly}'s two steps apply verbatim: the limit gives $A_PA_Q=1$, and subtracting
leaves $A_P\varepsilon_Q+A_Q\varepsilon_P+\varepsilon_P\varepsilon_Q=0$ with every term nonnegative
and some term strictly positive if a nonconstant member exists.
\end{proof}

The proof of Proposition~\ref{K-prop:nopoly} uses polynomiality in exactly one place, to pass from
``holding at infinitely many $t$'' to ``holding identically'', and restricting to the subsequence
removes that need. What the argument actually consumes is three hypotheses (finitely many members,
positive values, nonconstant members growing) and none of them is about shape.

The strengthening is not cosmetic. Excluding polynomial families rules out a reduction of \#307 to
Bunyakovsky, Schinzel or Bateman--Horn, but it says nothing about \emph{exponential} families, and
those are what the classical technique for problems of this kind actually uses: Th\=abit's rule for
amicable pairs runs on primes $3\cdot2^{n}-1$, $3\cdot2^{n-1}-1$, $9\cdot2^{2n-1}-1$, and Euler's
generalisation is of the same shape. A member $3\cdot2^{n}-1$ is positive and tends to infinity, so
it satisfies the hypotheses above and the family is excluded. Proposition~\ref{prop:nofamily}
therefore closes the whole Th\=abit/Euler programme for \#307, not merely its polynomial shadow: no
fixed--shape rule of any growth rate can generate solutions, and the constructive calculus of
Section~\ref{K-sec:frame} of the core note is not failing for want of a cleverer ansatz.

Two limits on the statement should be read with it. The hypothesis that the families have
\emph{fixed} cardinality is used, in the splitting into constants plus a remainder tending to $0$, and
a family whose support grows with $t$ is not covered; that gap is real, and harmless, since a growing
support is a growing number of simultaneous primality demands rather than a weaker one. And
positivity is a genuine hypothesis, not a consequence, for the reason recorded in
Remark~\ref{rem:nopolyerratum}. Formalised as \texttt{Erdos307.no\_family\_of\_any\_shape} on three
standard axioms, in the same shape as Proposition~\ref{K-prop:nopoly}: the abstract core is proved, and
the two facts the family supplies, that each nonconstant member contributes a strictly positive
reciprocal and that the remainders tend to $0$, enter as hypotheses rather than being derived, so
what is imported stays visible. The theorem's statement quantifies over arbitrary index types and
arbitrary $f_P:\mathbb{N}\to\iota\to\mathbb{R}$; no polynomial occurs anywhere in it.

\begin{remark}[Positivity is a hypothesis, not a consequence]\label{rem:nopolyerratum}
It is not enough to ask only that each polynomial take prime values for infinitely many integers $t$, without requiring the values to be positive on a common branch. Without positivity the statement is \emph{false}: take $P=(2,2,t,-t)$ and $Q=(2,2)$. Each member takes prime values for infinitely many
integers, and
\[ \Bigl(\tfrac12+\tfrac12+\tfrac1t-\tfrac1t\Bigr)\Bigl(\tfrac12+\tfrac12\Bigr)=1 \]
identically, with two nonconstant members present. The proof's step ``every term on the left is
nonnegative'' silently assumes positivity, which is exactly the missing hypothesis; with it, as now
stated, the proof is correct and the counterexample is excluded.

The formalisation carries the repaired statement, not the printed one:
\texttt{Erdos307.nonconstant\_empty} takes $0<f_P(i)$ for every nonconstant member as an explicit
hypothesis, so the Lean development was already stating the corrected proposition. This is the one
place in the paper where formalisation caught a missing hypothesis before review did.
\end{remark}

\noindent This upgrades the earlier computational observation (that the breeder family is not
polynomially parametrised) to a theorem for \emph{all} fixed--shape families: the prime--density
hypothesis behind \#307 is irreducibly non--polynomial, which is why no standard conjecture applies.

\begin{proposition}[Near--miss frames are extremal for $G$]\label{prop:nearframe}
Let $a$ be squarefree with $a'$ squarefree, and take the frame $(M_0,N)=(a,a')$, so that
$M_0'=a'$ and $N'=a''$. Then
\[ G\;=\;NM_0-N'M_0'\;=\;a'\,(a-a''), \]
and since $a''\ne a$ unless \#307 is solved, $|a-a''|\ge1$ and therefore $|G|\ge a'$.
\end{proposition}
\begin{proof}
Direct substitution; the inequality is $|a-a''|\ge1$ for the nonzero integer $a-a''$.
\end{proof}

The point is not that the bound is large but that it runs the wrong way. Writing $r=a''/a$ for the
near--miss ratio, $G=aa'(1-r)$, so a frame built from a good near--miss has $1-r$ small and $aa'$
correspondingly large, and the two cancel \emph{exactly}, because $a(1-r)=a-a''$ is an integer and a
nonzero one. Improving the near--miss therefore does not shrink $G$; it drives $a'$ up and $G$ with
it. The witness of Proposition~\ref{B-prop:nearmiss}, with $|1-r|=1.65\times10^{-13}$, gives
$|G|\ge a'>10^{112}$, against the $\log_{10}G\approx13$ that the Baker--Harman--Pintz route reaches
at the barrier. So the natural reading of ``steer $\Delta_0$ small by steering $s_{M_0}s_N$ to $1$''
is not merely circular in the sense of Section~\ref{K-sec:frame} of the core note: on near--miss frames it is inverted.

Small $|G|$ does occur, but only on small frames: over $1{,}143{,}888$ coprime squarefree frames with
$M_0,N\le2000$ the minimum is $|G|=5$, at $(M_0,N)=(2,3)$. That is the supply--demand conflict in raw
form, and it can be measured. Binning by $|G|$ and taking the best frame in each bin
(\texttt{code/breeder\_supply.gp}) gives
\[
\begin{array}{l|cccccc}
\log_{10}|G| & 0.5 & 1.5 & 2.5 & 3.5 & 4.5 & 5.0\\\hline
\max \tau(K)/|G| & 0.600 & 0.313 & 0.214 & 0.053 & 0.0065 & 0.0030
\end{array}
\]
The ratio is largest on the smallest frames, where it is already below $1$, and decays monotonically
from there. Since a productive frame needs $\tau(K)/|G|$ well above $1$, no frame in the range
supplies one, and the trend is against scale rather than with it.

\paragraph{The descent, and why it cannot be iterated.} Adjoining a prime $w$ to $M_0$ acts on
$G$ \emph{linearly}: with $M_0\mapsto M_0w$ one has $M_0'\mapsto M_0'w+M_0$ and hence
\[ G\ \longmapsto\ wG-M_0N', \qquad\text{so}\qquad |G_{\mathrm{new}}|=|G|\cdot|w-w^{*}|,
   \qquad w^{*}:=\frac{M_0N'}{G}=\frac{\sigma(N)}{\varepsilon}, \]
writing $\varepsilon=1-\sigma(M_0)\sigma(N)=G/(M_0N)$. Since $G=0$ is \#307, this is a descent
towards a solution, and on the $\varepsilon$ scale it is quadratic: $\varepsilon_{\mathrm{new}}
=\varepsilon^{2}\theta/\sigma(N)$ with $\theta=w-w^{*}$. It is a Newton iteration for \#307.

Three facts close it. First, it cannot start below the barrier: $w^{*}=\sigma(N)/\varepsilon\ge2$
needs $\varepsilon\le\sigma(N)/2$, i.e.\ $\sigma(M_0)\sigma(N)$ near $1$, which by AM--GM needs
$\sigma(M_0N)\ge2$ and so $\omega(M_0N)\ge59$. Over $4{,}284{,}896$ coprime squarefree frames with
$M_0\le5000$, $N\le3000$ the smallest $|\varepsilon|$ is $0.470$, and $w^{*}<2$ throughout.

Second, only $w=\lfloor w^{*}\rfloor$ or $\lceil w^{*}\rceil$ can shrink $|G|$, since
$|G_{\mathrm{new}}|<|G|$ demands $|w-w^{*}|<1$. So a step descends only when one of those two
integers is prime, with probability $\approx2/\log w^{*}$.

Third, and decisively, a successful step squares the difficulty. From $w^{*}=M_0N'/G$,
\[ w^{*}_{\mathrm{new}}=\frac{M_0wN'}{G\theta}=w^{*}\cdot\frac{w}{\theta}\ \approx\
   \frac{(w^{*})^{2}}{\theta}, \]
so $\log w^{*}$ more than doubles at every step while $\log|G|$ falls by only $|\log\theta|\approx
0.69$. The gain is linear in the number of steps $k$ and the cost is quadratic in the exponent:
$|G|$ falls like $2^{-k}$, whereas the probability of surviving $k$ steps is
$\prod_{j<k}2/\log w^{*}_j\asymp2^{-k^{2}/2}$. Starting from the greedy frame $N=30$,
$M_0=\prod_{7\le p\le271}p$, where $\log|G|=249$ and $w^{*}=428.2$
(\texttt{code/breeder\_descent.gp}):
\[
\begin{array}{l|cccccc}
\text{paths available} & 10^{2} & 10^{4} & 10^{6} & 10^{8} & 10^{10} & 10^{12}\\\hline
\text{reachable depth }k & 3 & 5 & 6 & 7 & 8 & 8\\
|G|\text{ shrinks by} & 8 & 32 & 64 & 128 & 256 & 256
\end{array}
\]
against the factor $10^{108}$ required to reach $G=0$. Even $10^{12}$ independent descent paths buy
eight steps and a factor $256$. This is the quantitative form of the circularity of
Section~\ref{K-sec:frame} of the core note: the Newton iteration for \#307 exists, is exact, and converges quadratically,
and each step costs a prime whose size is the square of the last one's.

\paragraph{Why $(\mathrm{H})$ is not a finite search.} It is tempting to read $(\mathrm{H})$ as
``run a certificate hunt that is expected to succeed.'' It is not, for a quantitative reason. The
expected number of certificates from one frame is $\mathbb{E}\approx
\bigl(\tau(K)/G\bigr)\cdot(\ln r\,\ln p\,\ln q)^{-1}$, and \emph{both} factors are pinned by the
geometry.
\begin{itemize}[leftmargin=1.4em]
\item \emph{Demand.} $G=\Delta_0=M_0N(1-s_{M_0}s_N)$ exactly, so $G$ is small only in the circular
regime $s_{M_0}s_N\approx1$; over $2\times10^4$ random coprime squarefree frames the ratio
$G/(M_0N)$ has median $1$ and best value $\approx10^{-0.23}$, and the honest Baker--Harman--Pintz
reduction still leaves $\log_{10}G\approx13$ at the barrier.
\item \emph{Supply.} $\tau(K)$ is bounded by the \emph{size} of $K=GN^2+c^2\ (\ge c^2)$, which is an
\emph{output} of the frame, not a free dial. The maximal order of the divisor function
\cite[Thm.~317]{HardyWright} caps $\log_2\tau(K)$ at $\approx36$ for $K\le10^{54}$ (the barrier
geometry) and at $\approx112$ even for $K\le10^{240}$; reaching $\tau(K)=2^{120}$ needs
$K\gtrsim10^{263}$, i.e.\ frame scale far beyond the barrier.
\end{itemize}
Enlarging a frame grows $K\propto N^2$ but gains divisors only at rate $\log2/\log\log K<1$ per
log--unit of $K$, while it grows $G$ at rate $1$ per log--unit of $N$; so $\tau(K)/G$ cannot be
driven up. Optimising $\mathbb{E}$ over all $(M_0,N)$ splits with geometry--consistent $K$ and the
maximal--order $\tau$, and granting the most favourable demand and singular--series constants, yields
$\log_{10}\mathbb{E}\le-4$ at the barrier and \emph{decreasing} with scale; empirically $385$
factorable frames realise a median of $0$ and a maximum of $1$ admissible divisor, against the
$\sim10^{5}$ a single productive frame would need. A self--consistent accounting; capping the
number $w$ of forced prime factors of $K$ by the frame's own residue freedom (counted generously),
sharpens this to $\log_{10}\mathbb{E}\le-8$ at $B=56$, falling to $\approx-260$ at $B=3000$: the
deficit is \emph{scale--invariant}, so no choice of scale, slot count, or budget split reaches
expectation $1$. (The construction is even self--similar: engineering $K$'s algebraic; e.g.\
Gaussian--integer; factorisation to supply divisors reduces to prescribing $N$ and $N'$
simultaneously, an arithmetic--antiderivative inverse problem of the same difficulty class as \#307.)
The breeder thus makes \#307 conditional on $(\mathrm{H})$ but does \emph{not} render it a feasible
finite computation: the supply--demand conflict through $K=GN^2+c^2$ is exactly the circularity of
Section~\ref{K-sec:frame} of the core note in quantitative form. There is also a structural reason the analogy with
amicable pairs was bound to fall short: te~Riele's breeding amplifies a \emph{known seed} pair into
new ones, whereas the barrier (Theorem~\ref{K-thm:barrier}) forbids any seed below $2.09\times10^{56}$
the method imports a precondition the problem cannot supply.


\begin{thebibliography}{9}
\bibitem{Barbeau1961} E.~J.~Barbeau, \emph{Remarks on an arithmetic derivative}, Canad.\ Math.\
Bull.\ \textbf{4} (1961), 117--122.
\bibitem{Barbeau1976} E.~J.~Barbeau, \emph{Computer challenge corner: Problem 477}, J.\ Recreational
Math.\ \textbf{9} (1976/77), 30.
\bibitem{ErdosGraham1980} P.~Erd\H{o}s and R.~L.~Graham, \emph{Old and New Problems and Results in
Combinatorial Number Theory}, Monogr.\ Enseign.\ Math.\ \textbf{28}, Geneva, 1980, p.~38.
\bibitem{GrahamEgyptian} R.~L.~Graham, \emph{Paul Erd\H{o}s and Egyptian fractions}, in \emph{Erd\H{o}s
Centennial}, Bolyai Soc.\ Math.\ Stud.\ \textbf{25}, Springer, 2013, pp.~289--309.
\bibitem{Watanabe2020} T.~Watanabe, \emph{New examples of the representation of $1$ by the sum of
reciprocals of semiprime numbers}, arXiv:2009.03275 (2020).
\bibitem{Bado2026b} I.-O.~Bado, \emph{Primary-pseudoperfect sectors in two-cycles of the
arithmetic derivative}, preprint (2026), communicated to the author.
\bibitem{BloomShifted} T.~F.~Bloom, \emph{Unit fractions with shifted prime denominators}, Proc.\
Roy.\ Soc.\ Edinburgh Sect.\ A (2024); arXiv:2305.02689.
\bibitem{UfnarovskiAhlander2003} V.~Ufnarovski and B.~\AA hlander, \emph{How to differentiate a
number}, J.\ Integer Seq.\ \textbf{6} (2003), Article 03.3.4.
\bibitem{Kovic2012} J.~Kovi\v{c}, \emph{The arithmetic derivative and antiderivative}, J.\ Integer
Seq.\ \textbf{15} (2012), Article 12.3.8.
\bibitem{MO320838} MathOverflow, \emph{Can the product of two sums of prime reciprocals be $1$?},
question 320838 (comments by R.~Israel and J.~Rosen, answer by ``Woett''), 2019.
\url{https://mathoverflow.net/questions/320838}.
\bibitem{ErdosProblems307} T.~F.~Bloom, \emph{Erd\H{o}s problem \#307},
\url{https://www.erdosproblems.com/307}.
\bibitem{OEIS_A003415} OEIS Foundation, \emph{Sequence A003415 (arithmetic derivative)},
\url{https://oeis.org/A003415}.
\bibitem{Pollack2018} P.~Pollack, \emph{On amicable tuples}, Illinois J. Math. \textbf{62} (2018),
  no.~1--4, 225--240; arXiv:1711.06847. (Explicit bounds for relatively prime amicable pairs,
  refining the finiteness of Artjuhov and Borho.)

\bibitem{teRiele1984} H.~J.~J.~te~Riele, \emph{Computation of all the amicable pairs below
$10^{10}$}, Math.\ Comp.\ \textbf{47} (1986), 361--368.
\bibitem{BHP2001} R.~C.~Baker, G.~Harman, and J.~Pintz, \emph{The difference between consecutive
primes, II}, Proc.\ London Math.\ Soc.\ (3) \textbf{83} (2001), 532--562.
\bibitem{Iwaniec1978} H.~Iwaniec, \emph{Almost-primes represented by quadratic polynomials},
Invent.\ Math.\ \textbf{47} (1978), 171--188.
\bibitem{LemkeOliver2012} R.~J.~Lemke Oliver, \emph{Almost-primes represented by quadratic
polynomials}, Acta Arith.\ \textbf{151} (2012), 241--261.
\bibitem{HardyWright} G.~H.~Hardy and E.~M.~Wright, \emph{An Introduction to the Theory of Numbers},
6th ed., Oxford Univ.\ Press, 2008.
\bibitem{Tenenbaum} G.~Tenenbaum, \emph{Introduction to Analytic and Probabilistic Number Theory},
3rd ed., Grad.\ Stud.\ Math.\ \textbf{163}, Amer.\ Math.\ Soc., 2015 (Erd\H{o}s--Wintner theorem).
\bibitem{Stothers} W.~W.~Stothers, \emph{Polynomial identities and Hauptmoduln}, Quart.\ J.\ Math.\
Oxford (2) \textbf{32} (1981), 349--370; R.~C.~Mason, \emph{Diophantine Equations over Function
Fields}, Cambridge Univ.\ Press, 1984.
\bibitem{Seva323194} Seva (\texttt{mathoverflow.net} user), \emph{A recursion with a number--theoretic
function}, MathOverflow question 323194 (2019). \url{https://mathoverflow.net/questions/323194}.
\bibitem{FewPrimes2025} \emph{Egyptian fractions for few primes}, arXiv:2507.03727 (July 2025).
\bibitem{Steinerberger2025} S.~Steinerberger, \emph{On an iterated arithmetic function problem of
Erd\H{o}s and Graham}, arXiv:2504.08023 (April 2025).
\bibitem{StretchedEG2026} \emph{A stretched--exponential bound for an Erd\H{o}s--Graham
unit--fraction problem}, arXiv:2607.04157 (July 2026).
\bibitem{PortFillings2026} \emph{Port fillings for primary pseudoperfect numbers}, arXiv:2605.21518
(May 2026).
\bibitem{Semiprime2026} \emph{Every natural number is a sum of distinct semiprime unit fractions},
arXiv:2606.15159 (June 2026).
\bibitem{ButskeJajeMayernik} W.~Butske, L.~M.~Jaje and D.~R.~Mayernik, \emph{On the equation
$\sum_{p\mid N}1/p+1/N=1$, pseudoperfect numbers, and perfectly weighted graphs}, Math.\ Comp.\
\textbf{69} (2000), 407--420.
\bibitem{Bado2026} I.~O.~Bado, \emph{Structural constraints for an Erd\H{o}s unit--fraction problem
over primes}, preprint, Aix--Marseille University (May 2026).
\url{https://doi.org/10.13140/RG.2.2.32245.54245}.
\bibitem{Giuga1950} G.~Giuga, \emph{Su una presumibile propriet\`a caratteristica dei numeri primi},
Ist.\ Lombardo Sci.\ Lett.\ Rend.\ Cl.\ Sci.\ Mat.\ Nat.\ (3) \textbf{14} (1950), 511--528.
\bibitem{BBBG1996} D.~Borwein, J.~M.~Borwein, P.~B.~Borwein, and R.~Girgensohn, \emph{Giuga's conjecture
on primality}, Amer.\ Math.\ Monthly \textbf{103} (1996), no.~1, 40--50.
\bibitem{BorweinWong1997} J.~M.~Borwein and E.~Wong, \emph{A survey of results relating to Giuga's
conjecture on primality}, CRM Proc.\ Lecture Notes \textbf{11} (1997), 13--27.
\bibitem{BJM2000} W.~Butske, L.~M.~Jaje, and D.~R.~Mayernik, \emph{On the equation
$\sum_{p\mid N}1/p+1/N=1$, pseudoperfect numbers, and perfectly weighted graphs}, Math.\ Comp.\
\textbf{69} (2000), no.~229, 407--420.
\bibitem{GrauOllerMarcen2012} J.~M.~Grau and A.~M.~Oller-Marc\'en, \emph{Giuga numbers and the
arithmetic derivative}, J.\ Integer Seq.\ \textbf{15} (2012), Article 12.4.1; arXiv:1103.2298.
\bibitem{SondowMacMillan2017} J.~Sondow and K.~MacMillan, \emph{Primary pseudoperfect numbers,
arithmetic progressions, and the Erd\H{o}s--Moser equation}, Amer.\ Math.\ Monthly \textbf{124} (2017),
no.~3, 232--240; arXiv:1812.06566.
\bibitem{GrauOllerMarcenSadornil2021} J.~M.~Grau, A.~M.~Oller-Marc\'en, and D.~Sadornil,
\emph{On $\mu$-Sondow numbers}, arXiv:2111.14211 (2021); Acta Math.\ Hungar.\ (2022).
\bibitem{Cox} D.~A.~Cox, \emph{Primes of the Form $x^2+ny^2$: Fermat, Class Field Theory, and Complex
Multiplication}, 2nd ed., Wiley, 2013.
\bibitem{OEIS_A054377} OEIS Foundation, \emph{Sequence A054377 (primary pseudoperfect numbers)},
\url{https://oeis.org/A054377}.
\bibitem{OEIS_A007850} OEIS Foundation, \emph{Sequence A007850 (Giuga numbers)},
\url{https://oeis.org/A007850}.
\bibitem{OEIS_A396915} OEIS Foundation, \emph{Sequence A396915 (composite squarefree $k$ with
$k'+2k$ a perfect square)}, \url{https://oeis.org/A396915}.
\bibitem{Evertse1984} J.-H.~Evertse, \emph{On sums of $S$-units and linear recurrences}, Compositio
Math.\ \textbf{53} (1984), 225--244.
\bibitem{FI1998} J.~Friedlander and H.~Iwaniec, \emph{The polynomial $X^2+Y^4$ captures its
primes}, Ann. of Math. \textbf{148} (1998), 945--1040.
\bibitem{HB2001} D.~R.~Heath-Brown, \emph{Primes represented by $x^3+2y^3$}, Acta Math.
\textbf{186} (2001), 1--84.
\bibitem{Stewart2013} C.~L.~Stewart, \emph{On divisors of Lucas and Lehmer numbers}, Acta Math.
\textbf{211} (2013), 291--314.
\bibitem{BCZ2003} Y.~Bugeaud, P.~Corvaja, and U.~Zannier, \emph{An upper bound for the G.C.D. of
$a^n-1$ and $b^n-1$}, Math. Z. \textbf{243} (2003), 79--84.
\bibitem{Merikoski2019} J.~Merikoski, \emph{On the largest prime factor of $n^2+1$}, J. Eur. Math.
Soc. \textbf{25} (2023), 1253--1284.
\bibitem{BGS2010} J.~Bourgain, A.~Gamburd, and P.~Sarnak, \emph{Affine linear sieve, expanders,
and sundry applications}, Invent. Math. \textbf{179} (2010), 559--644.
\bibitem{ESS2002} J.-H.~Evertse, H.~P.~Schlickewei, and W.~M.~Schmidt, \emph{Linear equations in
variables which lie in a multiplicative group}, Ann.\ of Math.\ (2) \textbf{155} (2002), 807--836.
\bibitem{Erdos1952} P.~Erd\H{o}s, \emph{On the sum $\sum_{k=1}^x d(f(k))$}, J.\ London Math.\ Soc.\
\textbf{27} (1952), 7--15.
\bibitem{BHV2001} Y.~Bilu, G.~Hanrot, and P.~M.~Voutier, \emph{Existence of primitive divisors of
Lucas and Lehmer numbers} (with an appendix by M.~Mignotte), J.\ Reine Angew.\ Math.\ \textbf{539}
(2001), 75--122.
\end{thebibliography}
\end{document}
